\documentclass{article}


% if you need to pass options to natbib, use, e.g.:
\PassOptionsToPackage{numbers, compress}{natbib}
% before loading neurips_2024


% ready for submission
%\usepackage{neurips_2024}


% to compile a preprint version, e.g., for submission to arXiv, add add the
% [preprint] option:
%\usepackage[preprint]{neurips_2024}


% to compile a camera-ready version, add the [final] option, e.g.:
\usepackage[final]{neurips_2024}


% to avoid loading the natbib package, add option nonatbib:
%    \usepackage[nonatbib]{neurips_2024}


\usepackage[utf8]{inputenc} % allow utf-8 input
\usepackage[T1]{fontenc}    % use 8-bit T1 fonts
\usepackage[hidelinks]{hyperref}      % hyperlinks
\usepackage{url}            % simple URL typesetting
\usepackage{booktabs}       % professional-quality tables
\usepackage{amsfonts}       % blackboard math symbols
\usepackage{nicefrac}       % compact symbols for 1/2, etc.
\usepackage{microtype}      % microtypography
%\usepackage{xcolor}         % colors
\usepackage[dvipsnames]{xcolor}
\usepackage{algorithm}
\usepackage{algorithmic}
\usepackage{microtype}
\usepackage{graphicx}
\usepackage{subfigure}

\usepackage{multirow}
\usepackage{multicol}
\usepackage{booktabs} % for professional tables
\usepackage{amsmath}
\usepackage{cleveref}
\usepackage{amssymb}
\usepackage{mathtools}
\usepackage{amsthm}
\usepackage{enumitem}
\usepackage{verbatim}
\usepackage{pifont}
%%%%%%%%%%%%%%%%%%%%%%%%%%%%%%%%
% THEOREMS
%%%%%%%%%%%%%%%%%%%%%%%%%%%%%%%%
\theoremstyle{plain}
\newtheorem{theorem}{Theorem}[section]
\newtheorem{proposition}[theorem]{Proposition}
\newtheorem{lemma}[theorem]{Lemma}
\newtheorem{corollary}[theorem]{Corollary}
\theoremstyle{definition}
\newtheorem{definition}[theorem]{Definition}
\newtheorem{assumption}[theorem]{Assumption}
\theoremstyle{remark}
\newtheorem{remark}[theorem]{Remark}

\newcommand{\simcom}[1]{\textcolor{blue}{[SM: #1]}}
\newcommand{\kbf}{\mathbf{k}}
\newcommand{\ebf}{\mathbf{e}}
\newcommand{\minuscule}{\fontsize{5}{6}\selectfont}


% \title{Fast and Complete G-Invariance via novel Algorithmic Reduction of the G-Bispectrum on Finite Groups}
\title{The Selective $G$-Bispectrum and its Inversion: Applications to $G$-Invariant Networks}


% The \author macro works with any number of authors. There are two commands
% used to separate the names and addresses of multiple authors: \And and \AND.
%
% Using \And between authors leaves it to LaTeX to determine where to break the
% lines. Using \AND forces a line break at that point. So, if LaTeX puts 3 of 4
% authors names on the first line, and the last on the second line, try using
% \AND instead of \And before the third author name.


\author{%
  Simon~Mataigne\thanks{Simon Mataigne is a Research Fellow of the Fonds de la Recherche Scientifique - FNRS.%Use footnote for providing further informationabout author (webpage, alternative address)---\emph{not} for acknowledgingfunding agencies.
  } \\
  ICTEAM, UCLouvain\\
  Louvain-la-Neuve, Belgium \\
  \texttt{simon.mataigne@uclouvain.be} \\
  % examples of more authors
   \And
   Johan Mathe\\
  Atmo\\
  San Francisco, CA\\
   \texttt{johan@atmo.ai} \\
  \AND
   Sophia Sanborn \\
  Science \\
  San Francisco, CA\\
\texttt{sophiasanborn@gmail.com} \\
  \And
  Christopher Hillar\\
  Algebraic\\
  San Francisco, CA\\
   \texttt{hillarmath@gmail.com} \\
   \And
   Nina Miolane\thanks{Nina Miolane acknowledges funding from the NSF grant 2313150.} \\
   UC Santa Barbara\\
  Santa Barbara, CA\\
   \texttt{ninamiolane@ucsb.edu} \\
}


\begin{document}


\maketitle


\begin{abstract}
% The development of Convolutional Neural Networks (CNNs) that are insensitive to a group action of the input, e.g., a rotation of the input image, is a critical task in machine learning.
% In Geometric Deep Learning \citep{bronstein2021geometric}, geometric priors derived from the structure of data are built into deep learning models to increase model efficiency and generalization. The layers of a Group-Equivariant Convolutional Neural Network (G-CNNs) \citep{cohenc16} are designed to have the property of \textit{equivariance} to a group of transformations, such as rotation -- that is, to

%Group-Equivariant Convolutional Neural Networks (G-CNNs) \citep{cohenc16} exploit the symmetries in data to reduce the translation-equivariance of convolutional neural networks can be generalized to other transformation groups

An important problem in signal processing and deep learning is to achieve \textit{invariance} to nuisance factors not relevant for the task. Since many of these factors are describable as the action of a group $G$ (e.g., rotations, translations, scalings), we want methods to be $G$-invariant. The $G$-Bispectrum extracts every characteristic of a given signal up to group action: for example, the shape of an object in an image, but not its orientation. Consequently, the $G$-Bispectrum has been incorporated into deep neural network architectures as a computational primitive for $G$-invariance\textemdash akin to a pooling mechanism, but with greater selectivity and robustness. However, the computational cost of the $G$-Bispectrum ($\mathcal{O}(|G|^2)$, with $|G|$ the size of the group) has limited its widespread adoption. Here, we show that the $G$-Bispectrum computation contains redundancies that can be reduced into a \textit{selective $G$-Bispectrum} with $\mathcal{O}(|G|)$ complexity. We prove desirable mathematical properties of the selective $G$-Bispectrum and demonstrate how its integration in neural networks enhances accuracy and robustness compared to traditional approaches, while enjoying considerable speeds-up compared to the full $G$-Bispectrum.

% Group-Equivariant Convolutional Neural Networks (G-CNNs) \citep{cohenc16, pmlr-v80-kondor18a} generalize the translation-equivariance of traditional CNNs to group-equivariance, using more general symmetry transformations such as rotation. For tasks such as classification, these transformations are removed at the end of the network to achieve group-invariance, typically by taking an average or maximum over the group. While this is indeed invariant, it is \textit{excessively} so; two inputs that are non-identical up to group action can yield the same output, resulting in a general lack of robustness to adversarial attacks \citep{sanborn2023bispectral, sanborn2023general}. The latter paper proposed an alternative method for achieving invariance without signal loss, called the $G$-triple correlation ($G$-TC). While this method yields demonstrable gains in accuracy and robustness, it comes with a substantial increase in computational cost. In this paper, we introduce a significantly more efficient invariant layer that scales only linearly with group size.  It is based on the $G$-\textit{Bispectrum}, which is the Fourier transform of the $G$-TC. Our main theoretical result is that for a large family of finite groups only $\mathcal{O}(|G|)$ coefficients of the $G$-Bispectrum are necessary to conserve the selective invariance of the $G$-TC.
% We validate this finding in a suite of computational experiments.

%% pre-May 6:
% Group-Equivariant Convolutional Neural Networks (G-CNNs) \citep{cohenc16} generalize the translation-equivariance of traditional CNNs to  \textit{group}-equivariance, using more general symmetry transformations such as rotation for weight tying.  For tasks such as classification, such transformations are removed at the end of the network, to achieve group-\textit{invariance}, typically by taking a maximum over the group. While this is indeed invariant, it is \textit{excessively} so; two inputs that are non-identical up to group action can yield the same output, resulting in a general lack of robustness to adversarial attacks \citep{sanborn2023bispectral, sanborn2023general}. The latter paper proposed an alternative method for achieving invariance without loss of signal structure, called the $G$-triple correlation ($G$-TC). While this method yields demonstrable gains in accuracy and robustness, it comes with a significant increase in computational cost. In this paper, we introduce a new invariant layer based on the Fourier transform of the $G$-TC: the $G$-Bispectrum. Operating in Fourier space allows us to significantly reduce the computational cost. Our main theoretical result provides a reduction of the $G$-Bispectrum that conserves the selective invariance of the $G$-TC, while only requiring $\mathcal{O}(|G|)$ coefficients. In a suite of experiments, we demonstrate that our approach retains the invariance of the $G$-TC, while significantly reducing the computational cost.
\end{abstract}

\section{Introduction}\label{sec:introduction}
The visual world is rich with symmetries. For example, the identity of an object is invariant to its position in the visual field; vision has \textit{translational symmetry}. Group theory is the mathematics used to describe transformations, their actions on objects, and the object's symmetry. As such, group theory has penetrated the fields of signal processing and deep learning alike. For example, the Fourier transform, pillar of signal processing, has been adapted to the $G$-Fourier transform, with its spectrum decomposing a signal defined over a group into several frequencies. More recently, researchers have become interested in the properties of higher-order spectra such as the \textit{Bispectrum}, and its generalization to signals over groups via the $G$-Bispectrum.

\paragraph{$G$-Bispectrum} The $G$-Bispectrum is the Fourier transform of the $G$-Triple Correlation ($G$-TC). Historically, higher-order spectra found initial applications in the context of classical signal processing as generalizations of the two-point autocorrelation \citep{tukey1953spectral, brillinger1991some, nikias1993signal}. The work of \citet{kakarala_thesis} illuminated the relevance of the $G$-Bispectrum for invariant theory, as it is the lowest-degree spectral invariant that is complete. Since then, it has appeared in diverse settings such as vision science \citep{zetzsche2001nonlinear}, machine learning \citep{kondor2007novel, kondor08}, and 3D modeling \citep{kakarala2012Bispectrum}.

% More recently, the group-invariant properties of this primitive have been explored by both signal processing and machine learning researchers.
%% previous:
% The work of \citet{kakarala_thesis} illuminated their relevance for invariant theory as the general triple correlation on \textit{groups} is the lowest-degree invariant polynomial that is complete \citep{sturmfels2008algorithms}.


\paragraph{Limitations of the $G$-Bispectrum for Deep Learning} The computational complexity of the $G$-Bispectrum has severely limited the reach of its applications. The most salient example of this limitation is in machine learning and deep learning. Convolutional Neural Network (CNN) \citep{lecun95convolutional,LeNet5} reflect and exploit the translational symmetry of the visual world. Group-Equivariant CNNs ($G$-CNNs) \citep{cohenc16, pmlr-v80-kondor18a} do just this, with more general group-equivariant convolutions to exploit %more general
symmetries like rotational symmetries. In both cases, one typically wants to preserve transformations throughout the layers of a network (i.e., to be group-\textit{equivariant}), and remove them only at the end when ``canonicalizing'' an image for classification (i.e., to be group-\textit{invariant}). While the theory of equivariant layers has been thoroughly developed \citep{cohen2021equivariant, weiler2023EquivariantAndCoordinateIndependentCNNs}, less attention has been paid to the theory of invariant layers \citep{higgins2018definition}. This is where the $G$-Bispectrum enters the picture, and where its computational cost has strongly limited its integration into deep learning.

Commonly, invariance in $G$-CNNs is achieved by simply taking an average or maximum over the transformation group (Average or Max $G$-Pooling, respectively). However, as noted by \citet{sanborn2023general}, this is a highly lossy operation removing information about the structure of the signal. While the \texttt{max} operation is indeed invariant (the \texttt{max} of an image is the same as the \texttt{max} of an image rotated by $90$ degrees), it is \textit{excessively} invariant: one could permute all of the pixels in the image and without changing the maximum, but with none of the same structure (see Figure~\ref{fig:sophias_figure}). To address this, \citet{sanborn2023general} used the $G$-TC as a $G$-invariant layer that is \textit{complete}\textemdash that is, it removes group transformations with no loss of signal structure. This approach achieves demonstrable gains in accuracy and robustness \citep{sanborn2023general}, but it is computationally expensive.

Indeed, the space complexity of the $G$-TC, i.e, its number of coefficients, scales as $\mathcal{O}(|G|^2)$, where $|G|$ is the size of the group. As each coefficient demands for $\mathcal{O}(|G|)$ operations, the computational cost or the time complexity of the $G$-TC is $\mathcal{O}(|G|^3)$. An alternative would be to use the $G$-Bispectrum as the pooling layer. However, both its space and time complexities are $\mathcal{O}(|G|^2)$. By comparison, the Max $G$-pooling layer features a $\mathcal{O}(|G|)$ computational cost and returns a scalar output. This raises the question of whether one can achieve complete invariance and adversarial robustness without sacrificing too much in terms of computational efficiency.

\textbf{Contributions.} In this work, we prove for the first time that we can significantly reduce the computational complexity of the $G$-Bispectrum. This result has important implications for signal processing and deep learning on groups, for which the $G$-Bispectrum is a foundational computational primitive. Our contributions are:
\begin{itemize}[leftmargin=*]
    \item We provide a general algorithm that reduces the computational complexity of the $G$-Bispectrum from $\mathcal{O}(|G|^2)$ to $\mathcal{O}(|G|)$ in space complexity and from $\mathcal{O}(|G|^2)$ to $\mathcal{O}(|G|\log |G|)$ in time complexity if an FFT is available on $G$. %We leverage symmetries within the $G$-Bispectrum and that make its computation redundant.
    We term it the \textit{selective $G$-Bispectrum}. The algorithm can be applied to any finite group.
    \item We prove that the selective $G$-Bispectrum is complete for the most important finite groups used in practice, i.e., all discrete commutative groups, the dihedral groups of any order, the octahedral and full octahedral group. This significantly extends the work of \citep{giannakis1989signal, kakarala_thesis, sadler1992shift}, who first showed this for \textit{some} finite, commutative groups, where it was demonstrated that the $G$-Bispectrum can be computed with only $|G|$ space complexity.
    \item We use the selective $G$-Bispectrum to propose a new \textit{$G$-invariant} layer that strikes a balance between robustness and efficiency. In particular, it is more expensive than the Max $G$-pooling, % which has $O(|G|)$ time and $O(1)$ space complexity,
    but cheaper than the $G$-TC pooling.%, which suffer from $O(|G|^3)$ time and $O(|G|^2)$ space complexity.
     It is also cheaper than the full $G$-bispectral pooling of $\mathcal{O}(|G|^2)$ time and $\mathcal{O}(|G|^2)$ space complexity. The selective $G$-Bispectrum is more robust than the max $G$-pooling, and almost as robust as the $G$-TC.
    \item We run extensive experiments on the MNIST \citep{lecun2010mnist} and EMNIST \citep{cohen2017emnist} % and CIFAR-10 \citep{CIFAR10}
    datasets to evaluate how each invariant layer (Max $G$-pooling, $G$-TC, selective $G$-Bispectrum) impacts accuracy and speed on classification tasks. We achieve the expected results: Our layer is faster than the $G$-TC and full $G$-Bispectrum and more accurate than Max $G$-pooling.
    %\item We demonstrate the adversarial robustness of our approach against a class of well-known adversarial attacks leveraging the concept of metamers. Our selective bispectral layer achieves the expected results: more robust than architectures using the Max $G$-pooling, and equivalently robust compared to the $G$-CNNs equipped with $G$-TC and full $G$-bispectral poolings.
    \item We present several findings important to the design of invariant layers to guide further advances in the field of geometric deep learning. In particular, we show that the accuracy and speed advantages of the selective $G$-Bispectrum is most striking for $G$-CNNs with low number of convolutional filters. Conversely, increasing the number of filters in the $G$-Convolutions allows the Max $G$-Pooling to catch up on the accuracy.% of the selective $G$-bispectral pooling.
     This demonstrates that the $G$-bispectral pooling will be particularly interesting for neural networks operating under a smaller parameter budget.
\end{itemize}
We hope that the proposed reduction of the $G$-Bispectrum complexity will further open areas of research in signal processing on groups, that were previously prohibited due to the high complexity of the operation.




%Our work thus brings the precise and mathematically-motivated findings of \citep{sanborn2023general} into the sphere of practicality through cost-efficiency.

% Finally, we  the choice of a reduced Bispectrum layer against a full Bispectrum layer.

% This allows us to introduce a novel layer for $G$-invariant pooling that strikes a balance between robustness and efficiency.

% The cost to compute a bispectral coefficient is usually independent of $|G|$ if one knows the Fourier transform of the input signal (of size $|G|$). For all finite groups $G$, the Fourier transform can be obtained in $\mathcal{O}(|G|^2)$ operations at most, even $\mathcal{O}(|G|\log(|G|))$ when a Fast Fourier Transform (FFT) algorithm is known on $G$. In particular, the Cooley-Tukey FFT algorithm can be extended to many finite groups, including the dihedral group \citep{Diaconis1990EfficientCO, michael89fft}. Therefore, we can reduce the cost of the Bispectrum calculation on many classes of groups to $\mathcal{O}(|G|\log|G|)$ in computation and $\mathcal{O}(|G|)$ in space.

% This motivates the introduction of a $\mathcal{O}(|G|^2)$ layer in both space and computations (or less). This layer should provide more accurate results than the Max $G$-pooling as well as being less expensive than the $G$-TC.
% The improvement comes from the following observation: there are many redundancies in the $G$-TC that can be neatly exploited in the Fourier domain (i.e. with the \textit{Bispectrum} \citep{kakarala_finite_groups}), permitting a substantial reduction in size and cost to compute.



% This is most readily seen in the Fourier domain -- i.e. with the \textit{Bispectrum} \citep{kakarala_finite_groups}.



% Moreover, we can design a layer returning the full Bispectrum, in which case, the complexity is $\mathcal{O}(|G|^2)$ in both computations and space. To determine which coefficients induce completeness, it is enough to determine which allow to reconstruct the initial signal entirely. Retrieving the signal from bispectral coefficients is called the \emph{Bispectrum inversion problem} \citep{kakarala_thesis,Bendory_2018}.

% In this paper, we propose an implementation of the $G$-CNN where the pooling layer is a Bispectrum layer, as suggested in \citep{sanborn2023general}. This corresponds to switching the neural network in the Fourier domain. We will perform experiments on both the selective Bispectrum layer (only a few coefficients) and the full Bispectrum layer (all coefficients). To ensure completeness of the selective Bispectrum, we show that the Bispectrum can be inverted using $|G|$ bispectral coefficients for all finite commutative groups and $\frac{|G|}{4}$ coefficients for the dihedral group, i.e., the symmetries of the $n$-gon. The dihedral group is the main non-commutative group of interest for 2D-image processing.\simcom{Nina, a reference from clinical application or so?}.

% Extending the implementation og \citet{sanborn2023bispectral} to assess the performance of the reduced bispectra layers on the MNIST \citep{deng2012mnist} and EMNIST \citep{cohen2017emnist} datasets.
% We show that the bispectral layer achieves the expected results: it is both more accurate than the Max $G$-pooling and cheaper than the $G$-TC. Finally, we discuss the choice of a reduced Bispectrum layer against a full Bispectrum layer.

\section{Background: $G$-Triple Correlation and $G$-Bispectrum}\label{sec:triple_corr_and_bsp}

The proposed selective $G$-Bispectrum operation is closely related to two other foundational operations on signals defined on groups: the $G$-Triple Correlation and the full $G$-Bispectrum, which we introduce here. The background on group theory, including the definitions of groups, group actions, equivariance and invariance, is presented in Appendix~\ref{app:group_theory}.

% We need first to present alternative $G$-pooling layers that have recently emerged in addition to the well-known Max $G$-Pooling. Specifically, \citep{sanborn2023bispectral} uses the \textit{full} $G$-Bispectrum layer and \citep{sanborn2023general} introduces the $G$-triple correlation ($G$-TC) layer. We will introduce our \textit{selective} $G$-Bispectrum layer and provide mathematical results that prove its computational efficiency compared to the $G$-TC.

\paragraph{The $G$-Triple Correlation}
%The $G$-triple correlation ($G$-TC) \citep{kakarala_thesis} has had many uses but our focus stems from its utilization in $G$-CNNs.
Given a real signal defined on a finite group $\Theta:G\mapsto\mathbb{R}$, the $G$-Triple Correlation ($G$-TC)~\citep{kakarala_thesis}  is the lowest order polynomial that is complete, i.e., that conserves all of the information of the signal $\Theta$, up to group action by $G$.
\begin{definition}\label{def:gtc}
    The \emph{$G$-Triple Correlation} of a real signal $\Theta:G\mapsto\mathbb{R}$ is given by
    \begin{equation}
        T(\Theta)_{g_1,g_2}:=\sum_{g\in G} \Theta(g)\Theta(g\cdot g_1)\Theta(g\cdot g_2) \text{ for all } g_1, g_2\in G.
    \end{equation}
\end{definition}

The original triple-correlation was introduced for the classical framework of translations of a one-dimensional signal, i.e., where $X =\mathbb{Z}$ and $(G,\ \cdot) = (\mathbb{Z},+)$. The $G$-triple correlation from Definition~\ref{def:gtc} extends the original definition to any finite group $(G,\ \cdot)$. In our setting, the signal $\Theta:G\mapsto\mathbb{R}$ will be obtained after the $G$-convolution of a function $f:X\mapsto\mathbb{R}^c$,
%(e.g., a continuous RGB picture is a function $f:\mathbb{R}^2\mapsto\mathbb{R}^3$ function)
representing a continuous image with $c$ channels, with a filter $\phi:X\mapsto\mathbb{R}^c$, and the $G$-TC will be applied channel-by-channel. Importantly, the $G$-TC layer has computational complexity $\mathcal{O}(|G|^3)$ and outputs $\mathcal{O}(|G|^2)$ coefficients.

%For all $g \in G$, we have
% \begin{equation}
%     \Theta(g):=(\phi*f)(g) = \sum_{x\in X}\phi(\alpha(g^{-1},x))^{\top}f(x).
% \end{equation}
%  If $X =\mathbb{Z}$ and $(G,\ \cdot) = (\mathbb{Z},+)$, we obtain the classical discrete convolution.

\paragraph{The $G$-Bispectrum} The $G$-TC operation has a Fourier equivalent: the $G$-Bispectrum. Indeed, the definition of the Discrete Fourier Transform (DFT) can be extended to any finite group~(see, e.g., \citep{Diaconis1990EfficientCO}), as recalled below.
\begin{definition}
    Given a set of unitary representatives (Def.~\ref{def:representation}) of the equivalence classes of irreps $\rho_i:G\mapsto\mathrm{GL}(V_i)$ (Def.~\ref{def:equivalence}),
    the \emph{$G$-Fourier Transform} on a finite group $G$ of a signal $\Theta:G\mapsto\mathbb{R}$  is defined as
    \begin{equation}
        \mathcal{F}(\Theta)_{\rho_i} :=\sum_{g\in G}\Theta(g)\rho_i(g)^\dagger,
    \end{equation}
     where $\rho_i(g)^\dagger$ refers to the conjugate transpose of the matrix $\rho_i(g)$  (or simply transpose if $\rho_i$ is real-valued).
\end{definition}

The $G$-Bispectrum $\beta(\Theta)$ is defined as $\mathcal{F}(T(\Theta))$, with $\mathcal{F}$ evaluated over the group $G\times G$.
\citet{kakarala_thesis} proposed a closed-form expression for the $G$-Bispectrum $\beta(\Theta)$ directly in terms of $\mathcal{F}(\Theta)$. We recall it in Theorem~\ref{def:Bispectrum}.

\begin{theorem}{\citep{kakarala_finite_groups}}\label{def:Bispectrum}
    The $G$-Bispectrum  of a signal $\Theta:G\mapsto\mathbb{R}$, $\beta(\Theta)$, is given by:
    \begin{equation*}
        \beta(\Theta)_{\rho_1,\rho_2} =\left[\mathcal{F}(\Theta)_{\rho_1}\otimes\mathcal{F}(\Theta)_{\rho_2}\right]C_{\rho_1,\rho_2}\left[\bigoplus_{\rho\in\rho_1\otimes\rho_2}\mathcal{F}(\Theta)_\rho^\dagger\right]C_{\rho_1,\rho_2}^{\dagger},
    \end{equation*}
    where $C_{\rho_1,\rho_2}$ is a unitary matrix called the Clebsch-Gordan matrix, whose definition is recalled in Appendix~\ref{app:group_theory}.
    For each pair $\rho_1, \rho_2$, the matrix $\beta(\Theta)_{\rho_1, \rho_2}$ is called a $G$-bispectral coefficient.
\end{theorem}

For commutative groups, Theorem~\ref{def:Bispectrum} simplifies to a more compact expression, recalled in Theorem~\ref{thm:commutative_Bispectrum}. However, both the space and time complexity of the $G$-Bispectrum remain $O(|G|^2)$.

To the authors' best knowledge, there is no generic analytical formula for computing $C_{\rho_1,\rho_2}$ for an arbitrary group $G$. However, there exist formulas for specific classes of groups (see Appendix~\ref{app:clebsch_gordan}). Additionally, there exist packages for computing these for many groups using packages such as \texttt{escnn} \citep{cesa2022a}.

\paragraph{Complete $G$-Invariants} The $G$-TC and the $G$-Bispectrum are desirable computational primitives for signal processing and deep learning because they are \textit{complete $G$-invariants} (for generic data $\Theta, \widetilde{\Theta}$). Indeed, this completeness property make them very interesting for building invariant layers in $G$-CNNs, as they are selectively invariant. We define \textit{complete $G$-invariance} next.
\begin{theorem}{\citep[Thm.3.2] {Kakarala2009CompletenessOB}}
    \label{thm:bsp_invariance}
    The $G$-TC and the $G$-Bispectrum are complete $G$-invariants, i.e., for $\Theta, \widetilde{\Theta}:G\mapsto\mathbb{R}$ with $\mathcal{F}(\Theta)_\rho$ nonsingular for all irreps $\rho$, $T(\Theta)=T(\widetilde{\Theta})$, respectively  $\beta(\Theta) = \beta(\widetilde{\Theta})$, if and only if there exists $h\in G$ such that $\Theta(g) = \alpha(h, \widetilde{\Theta}(g))$ for all $g\in G$.
\end{theorem}
\paragraph{Application: $G$-invariant layers}\label{sec:bispectralgcnn}
\begin{figure*}
    \centering
    \includegraphics[width = 14cm]{Images/figures_ICML.pdf}
    \caption{Illustration of the different proposed $G$-CNN modules \citep{cohenc16, sanborn2023general}. The input $f$ is first processed through the $G$-convolutional layer composed of $K$ filters $\{\phi_k\}_{k=1}^K$. Then, an invariant layer is chosen (Max $G$-pooling, $G$-TC, or the selective/full $G$-Bispectrum layer). Finally, the ``pooled'' output is fed to a neural network designed for the machine learning task at hand. \vspace{-0.5cm}}
    \label{fig:GCNN}
\end{figure*}
The $G$-CNN architecture, first proposed in \citep{cohenc16}, is illustrated in Figure~\ref{fig:GCNN}. The input signal $f : X\mapsto\mathbb{R}$, typically an image, is processed through a $G$-Convolution layer using filters $\{\phi_k\}_{k=1}^K$. The output is feature maps $\{\Theta_k\}_{k=1}^K$ that form a set of $K$ real-valued signals with domain $G$. This $G$-Convolution layer is traditionally followed by a $G$-invariant layer. The most common is the Max $G$-Pooling layer. More recent works have proposed two alternatives based on the $G$-TC and the full $G$-Bispectrum: the $G$-TC Pooling~\citep{sanborn2023general} and the (full) $G$-Bispectrum~\citep{sanborn2023bispectral} respectively, where the latter requires the computations of the Fourier transforms of the feature maps, preferably computed using a Fast Fourier Transform (FFT) algorithm on $G$ \citep{Diaconis1990EfficientCO}. When testing the impact of the choice of $G$-invariant layer, the output of the invariant layer is typically fed to a Secondary Neural Network (NN) to perform the desired task, e.g., image classification. The Secondary NN often takes the form of a Multi-Layer Perceptron (MLP).

Experimental results have demonstrated the superior accuracy and adversarial robustness of the $G$-CNN equipped with a $G$-TC and $G$-Bispectrum invariant layer~\citep{sanborn2023bispectral,sanborn2023general}. However, both methods inherit the high space and time complexity of their respective operations. % space and time complexity of the $G$-TC are $K\mathcal{O}(|G|^2)$ and $\mathcal{O}(K|G|^3)$, while space and time complexity of the full $G$-Bispectrum are both $\mathcal{O}(K|G|^2)$ where $K$ is the number of filters of the $G$-Convolution.
This raises the question of whether we can reduce this computational complexity.



%\section{Bispectrum Inversion for Complete $G$-Invariant Neural Networks}\label{sec:bspinversion}
\section{Method: The Selective $G$-Bispectrum and its Inversion}\label{sec:bspinversion}

\paragraph{The Selective $G$-Bispectrum} We introduce a novel tool for signal processing on groups: the selective $G$-Bispectrum. %We obtain this object by introducing a novel derivation of the full $G$-Bispectrum from a subset of its coefficients.
A selective $G$-Bispectrum $\beta_{sel}$ is any $\mathcal{O}(|G|)$  subset of all coefficients of the $G$-Bispectrum $\beta$ (Definition~\ref{def:Bispectrum}), only conserving well-chosen pairs of irreps $(\rho_1, \rho_2)$. Which pairs of irreps to select depends on the group of interest.  This is possible due to redundancies and symmetries in the full object. Below, we provide an algorithmic procedure to compute the selective $G$-Bispectrum for any finite group $G$ that features at most $|\text{Irreps}|(\leq |G|)$ coefficients. The procedure is summarized in Algorithm \ref{alg:selective_bsp}. We have the following proposition.

\begin{algorithm}
    \caption{Selective $G$-Bispectrum on any finite group $G$}\label{alg:selective_bsp}
    \small
    \begin{algorithmic}[1]
        \STATE \textbf{Input}: Signal $\Theta: G \mapsto \mathbb{R}$ with $n = |G|$. Empty list of coefficients $L_\beta$. Empty list of irreps $L_\rho$. Kronecker Table of $G$.
        \STATE Add $\beta(\Theta)_{\rho_0,\rho_0}$ to $L_\beta$ where $\rho_0$ is the trivial irreps ($\rho_0(g)=1$ for all $g\in G$).
        \STATE Add $\rho_0$ to $L_\rho$.
        \STATE Choose $\tilde \rho_1$ such that $\tilde \rho_1 \otimes\tilde  \rho_1 = C_{\tilde \rho_1,\tilde \rho_1}\left(\bigoplus_{\rho \in \mathcal{R}} \rho\right) C_{\tilde \rho_1,\tilde \rho_1}^\dagger$ generates at least one irreps $\rho$ not yet in $L_\rho$.
        \STATE Add $\beta(\Theta)_{\rho_0,\tilde \rho_1}$ and $\beta(\Theta)_{\tilde \rho_1,\tilde \rho_1}$ to $L_\beta$, add $\tilde \rho_1$ to $L_\rho$.
        \STATE Add every $\rho$ that appears in $\tilde \rho_1 \otimes\tilde  \rho_1$ to $L_\rho$.
        \WHILE{$L_\rho$ keeps changing:}
            \STATE Find $\rho', \rho''$ in $L_\rho$ such that $\rho' \otimes \rho'' $ generates at least one irreps not already in $L_\rho$.
            \STATE Add $\beta(\Theta)_{\rho',\rho''}$ to $L_\beta$.
            \STATE Add $\rho$ that appears in $\rho' \otimes \rho''$ to $L_\rho$.
        \ENDWHILE
        \STATE \textbf{Return} $\beta_{sel}(\Theta):=L_\beta$: the selected $G$-bispectral coefficients.
    \end{algorithmic}
\end{algorithm}

\begin{proposition}
    The selective $G$-Bispectrum $\beta_{sel}$ from Algorithm~\ref{alg:selective_bsp} has at most $|G|$ coefficients.
\end{proposition}
\begin{proof}
    By construction of Algorithm \ref{alg:selective_bsp}, $|L_\rho|\leq |\text{Irreps}|$. Since $|\text{Irreps}|$ is at most the number of conjugacy classes of $G$ (see, e.g.,  \citet[Corollary~4.3.10]{steinberg2011representation}), we have $|\text{Irreps}|\leq|G|$.
\end{proof}
In Algorithm~\ref{alg:selective_bsp}, we note that the choice of $\rho_1$ is important,  since some $\rho_1$ will not allow the user to recover all of the irreps and therefore not ensure the completeness of the selective. We illustrate the computation of the selective $G$-Bispectrum in Figure~\ref{fig:selective_Bispectrum}, where we choose $\tilde \rho_1 = \rho_6$.
\begin{figure}
    \centering
    \includegraphics[width=1\linewidth]{Images/selective_Bispectrum_algorithm-01.png}
    \vspace{-0.5cm}
    \caption{Computation of the selective $G$-Bispectrum for the Full Octahedral Group. The gradient of color represents the order in which the $G$-bispectral coefficients are computed. The Kronecker Table represents which irreps emerge from the decomposition into irreps of the tensor product $\rho_i \otimes \rho_j$. We observe that the selective $G$-Bispectrum has only $6$ coefficients, compared to $100$ coefficients for the full $G$-Bispectrum.}
    \label{fig:selective_Bispectrum}
\end{figure}
\paragraph{Inverting the Selective $G$-Bispectrum for completeness} The \emph{inversion} of the selective $G$-Bispectrum $\beta_{sel}(\Theta)$ is reconstructing a signal $\mathcal{F}(\widetilde{\Theta})$ from the $G$-Bispectrum coefficients in the list $L_\beta$ such that $\widetilde{\Theta}=\alpha(g, \Theta)$ for some $g\in G$ (The $G$-Bispectrum is $G$-invariant, hence, $\Theta$ can only be recovered at best up to group action). Once $\mathcal{F}(\widetilde{\Theta})$ is known, $\widetilde{\Theta}$ can be obtained using the Inverse Fourier Transform.
If the selective $G$-Bispectrum can be inverted, then, by definition, it is complete in the sense of Theorem~\ref{thm:bsp_invariance}.
%If one c $\tilde \Theta$ is equal to the original signal $\Theta$ up to group action. Importantly, both the selective $G$-Bispectrum and its inversion can be computed for every finite group.


\section{Theory: Completeness of the Selective $G$-Bispectrum}\label{sec:completeness}

Our main theoretical claim is that the selective $G$-Bispectrum can be inverted and is a complete $G$-invariant that drastically reduces the complexity of the $G$-Bispectrum. We prove this claim for many finite groups $G$ of interest in signal processing and deep learning in a sequence of theorems presented in this section.

%The main idea behind the proofs is the previously mentioned \emph{$G$-Bispectrum inversion problem}. For a given group $(G,\cdot)$, we prove that the $G$-Bispectrum can be inverted and is hence complete. % Specifically, for each finite group $G$, we prove that the selective $G$-Bispectrum $\tilde \beta(\Theta)$ is enough to recover all entirely $\mathcal{F}(\Theta)_\rho$, for all of the irreps $\rho$. This means that the selective $G$-Bispectrum is enough to recover the original signal $\Theta$ up to group action. This, in turn, means that the selective $G$-Bispectrum is enough to recover the full $G$-Bispectrum.

\paragraph{Known Theorems} Previous authors had looked into the $G$-Bispectrum inversion problem. It is well known that $|G|=n$ coefficients are enough for the cyclic group $(C_n, \cdot) = (\mathbb{Z} / \mathbb{Z}_n, +\mod n)$.

\begin{theorem}{\citep{kakarala_finite_groups}}\label{thm:cyclic_inversion}
For cyclic groups $C_n$, $n\in\mathbb{N}_0$, the $C_n$-Bispectrum can be inverted using $|G|=n$ coefficients if $\mathcal{F}(\Theta)_\rho\neq 0$ for all irreps $\rho$ of $C_n$.
\end{theorem}

Similarly for a product of two such groups, we have the following theorem.

\begin{theorem}{\citep{giannakis1989signal}}\label{thm:product_cyclic_inversion}
For a product of cyclic groups $C_n \times C_m$, $n,m\in\mathbb{N}_0$, the $G$-Bispectrum can be inverted using $|G|=n m$ coefficients if $\mathcal{F}(\Theta)_\rho\neq 0$ for all $\rho \in C_n \times C_m$.
\end{theorem}

\paragraph{New Theorems} From now on, we assume that the Fourier transform $\mathcal{F}(\Theta)$ only features \textit{non-zero elements}, or \textit{invertible matrices} in the case of non-scalar Fourier coefficients. This assumption is supported by the zero probability of encountering this corner case (an arbitrarily small perturbation of any signal makes this assumption true).

We first extend the above results to all commutative groups. The proof relies on the fact that every finite commutative group is the direct sum of finitely many cyclic groups.

\begin{theorem}
For finite commutative groups $G$, the $G$-Bispectrum can be inverted using $|G|$ coefficients if $\mathcal{F}(\Theta)_\rho \neq 0$ for all $\rho \in G$.
\end{theorem}
See Appendix~\ref{app:commutative_inversion} for the proof and derivation of the inversion for the specific case of commutative groups. We note that our approach to inversion is symbolic, in that a solution can be expressed explicitly as a formula in terms of the input.  Other approaches are also possible to determine an inverse, such as using least squares \citep{haniff1991least} or more recent spectral methods \citep{chen2018spectral}.

We now extend the result to dihedral groups. Dihedral groups are ubiquitous in signal processing and deep learning because they represent the group of rotations and reflections.

\begin{theorem}
    For any dihedral group $D_n$ (symmetries of the $n$-gon), $n\in\mathbb{N}_0$, we need at most $\left \lfloor{ \frac{n-1}{2}}\right \rfloor + 2$ bispectral matrix coefficients for inversion if $\det(\mathcal{F}(\Theta)_\rho) \neq 0$ for all irreps $\rho$ of $D_n$. This corresponds to $1+4+16\cdot\left\lfloor{ \frac{n-1}{2}}\right \rfloor\approx 4 |D_n|$ scalar values.
\end{theorem}
The proof is provided in Appendix~\ref{app:dihedral_inversion}.
We now extend the result to octahedral and full octahedral groups, that are related to the symmetries of the octahedron. These groups are very important in signal processing and deep learning of 3D images.

\begin{theorem}
    For the octahedral group $O$ which has $|G| = 24$ group elements and $5$ irreps, we need only $4$ $G$-Bispectral coefficients in the selective $G$-Bispectrum. this corresponds to $172$ scalars. For the full octahedral group $FO$ which has $|G| = 48$ elements, we only need $6$ $G$-Bispectral coefficients in the selective $G$-Bispectrum to perform inversion. This corresponds to $334$ scalars.
\end{theorem}

A sketch of proof is provided in Appendix~\ref{app:octahedral} given the redundancy of the procedure. We see that the selective $G$-Bispectrum uses only $4$ coefficients, compared to $25$ coefficients needed for the full $G$-Bispectrum of the octahedral group. For the full octahedral group, it requires only $6$ coefficients compared to the $100$ coefficients of the full $G$-Bispectrum. In Figure~\ref{fig:overview}, we compare the full and selective $G$-Bispectra of the dihedral group $D_4$ (symmetries of the square) and the octahedral group.

\begin{figure}
    \centering
    \includegraphics[width=1\linewidth]{Images/selective_bispectrum_overview.png}
    \vspace{-0.5cm}
    \caption{Comparison of full and selective $G$-Bispectra for the dihedral group $D_4$ (left) and the octahedral group $O_h$ (right). The Kronecker tables of both groups show which irreps emerge from the decomposition into irreps of the tensor product of irreps $\rho_i \otimes \rho_j$. The colored boxes highlight the bispectral coefficients chosen for the full and selective Bispectra. Our proposed selective Bispectrum captures the same information as the full Bispectrum but with significantly fewer coefficients.}
    \label{fig:overview}
\end{figure}
\section{Experimental results}\label{sec:experiments}
\paragraph{Implementation and architecture} Our implementation of the selective $G$-bispectrum layer is based on the \href{https://github.com/gtc-invariance/gtc-invariance}{\texttt{gtc-invariance}} repository, implementing the $G$-CNN with $G$-convolution and $G$-TC layer \citep{sanborn2023general} and relying itself on the \href{https://github.com/QUVA-Lab/escnn}{\texttt{escnn} library} \citep{cesa2022a, weiler2021general}. The implementations related to this section can be found at the \href{https://github.com/geometric-intelligence/g-invariance}{\texttt{g-invariance}} repository.
%We propose a performance analysis of our method. To do so, we have implemented the Bispectrum layer for the cyclic group $C_n$ and the dihedral group $D_n$. These groups can be considered as discrete cases of the groups $\mathrm{SO}(2)$ (2D rotations) and $\mathrm{O}(2)$ (2D rotations and reflections) respectively.

We propose an experimental assessment of the newly proposed selective $G$-Bispectrum layer by comparing it with the Avg $G$-pooling, the Max $G$-pooling, the $G$-TC as invariance operations after the $G$-convolution of a $G$-CNN on the classification problems of the \href{http://yann.lecun.com/exdb/mnist/}{MNIST} dataset of handwritten digits \citep{lecun2010mnist}, the \href{https://www.nist.gov/itl/products-and-services/emnist-dataset}{EMNIST} dataset of handwritten letters \citep{cohen2017emnist} with standard train-test division. These datasets count 10 and 26 classes, respectively. We obtain transformed versions of the datasets -- $G$-MNIST/EMNIST -- by applying a random action $g\in G$ on each image in the original dataset.

The objective of our experiments is to isolate the speed-up of the $G$-Bispectrum layer. Hence, we consider architectures that only differ by the invariant layer in the classification task, following the experimental set up by \citep{sanborn2023general}.  The neural network architecture is composed of a $G$-convolution, a $G$-invariant layer, and finally a Multi-Layer-Perceptron (MLP), itself composed of three fully connected layers with ReLU nonlinearity. Finally, a fully connected linear layer is added to perform classification. The MLP's widths are tuned to match the number of parameters across each neural network model. The details are given in Appendix~\ref{app:training}. We highlight here that the pursued objective is to compare the differences in performances of the $G$-invariant layers, not to provide the state-of-the-art accuracy on the datasets involved. Henceforth, we do not optimize the architectures to reach the highest possible accuracy. We set simple architectures providing interpretable results for analysis. The experiments a performed using $8$ cores of a NVIDIA A30 GPU.
\paragraph{Training speed performance} Table~\ref{tab:complexities} recalls the theoretical complexities of the different layers. The computational cost of computing the selective $G$-Bispectrum is $\mathcal{O}(|G|\log |G|)$ if an FFT algorithm is available on $G$~\citep{Diaconis1990EfficientCO}, and $\mathcal{O}(|G|^2)$ with classical DFT. in Figure~\ref{fig:trainings}, we report the average training times on $\mathrm{SO}(2)/\mathrm{O}(2)$-MNIST for 10 runs as the discretization $C_n/D_n$ of $\mathrm{SO}(2)/\mathrm{O}(2)$ varies. In the first case, we use the FFT and observe that the Max $G$-pooling and $G$-Bispectrum training time scale linearly whereas it scales quadratically for the $G$-TC. For $\mathrm{O}(2)$, we perform a classic DFT on $D_n$ so that the $G$-Bispectrum scales worth. However, an FFT could be implemented to speed-up the process.
\begin{table}
    \centering
    \begin{tabular}{||c||c|c|c||}
        \hline
         Invariance layer&  Computational Complexity & Ouput size& Complete G-invariant\\
         \hline
         $G$-TC& $K\mathcal{O}(|G|^3)$&$K\mathcal{O}(|G|^2)$&\textcolor{Green}{\ding{51}}\\
         \hline
         Full $G$-Bispectrum& $K\mathcal{O}(|G|^2)$&$K\mathcal{O}(|G|^2)$&\textcolor{Green}{\ding{51}}\\
         \hline
         Select. $G$-Bispectrum & $K\mathcal{O}(|G|\log|G|\text{ or }|G|^2)$&$K\mathcal{O}(|G|)$&\textcolor{Green}{\ding{51}}\\
         \hline
         Max $G$-pooling&$K\mathcal{O}(|G|)$&$K\mathcal{O}(1)$&\textcolor{red}{\ding{55}}\\
         \hline
         Avg $G$-pooling&$K\mathcal{O}(|G|)$&$K\mathcal{O}(1)$&\textcolor{red}{\ding{55}}\\
         \hline
    \end{tabular}
    \vspace{0.2cm}
    \caption{$G$-CNN invariant layers and their computational cost and output size. $K$ is the number of filters. The selective $G$-Bispectrum that we propose is the complete $G$-invariant layer with the lowest time and space complexity. It reduces significantly the cost compared to the $G$-TC layer while preserving its completeness. \vspace{-0.4cm}}
    \label{tab:complexities}
\end{table}
\begin{figure}[h]
    \centering
    \includegraphics[width = 6.8cm]{Images/training_times_cyclic.pdf}
    \includegraphics[width =6.8cm]{Images/training_times_dihedral.pdf}
    \vspace{-0.3cm}
    \caption{Evolution of the average training times for the different invariant layers. The parameter $n$ is the size of the groups $C_n$ and $D_n$. The average and standard deviations are obtained over $10$ runs.  For all runs, the number of parameters of the complete neural network (filters and MLP) is set to $50000$ and $150000$ for $\mathrm{SO}(2)$ and $\mathrm{O}(2)$ respectively. Standard deviations are reported by vertical intervals. When a FFT is available, our selective G-Bispectrum significantly outperforms other complete G-invariant pooling layers in terms of speed. Specifically, when working with $C_{2^7}$, training on a dataset of $60000$ images takes only $247$ seconds, whereas the $G$-TC requires 1465 seconds. %As a result, the selective $G$-Bispectrum paves the way for the development of complete invariant pooling layers that can accommodate larger group sizes and, hence, a larger set of symmetries.
    }
    \label{fig:trainings}
\end{figure}

\paragraph{Classification Performance} We compare the performances of the $G$-Bispectrum layer with respect to the $G$-TC, the Max $G$-pooling and the Avg $G$-pooling models, trained on the $\mathrm{SO}(2)$/$\mathrm{O}(2)$-MNIST/EMNIST datasets and we assess the accuracy by averaging the validation accuracy over 10 runs. The classification accuracy is provided in Table~\ref{tab:classification}. For the experiments in Table \ref{tab:classification}, the following pattern holds: at equivalent number of parameters, the more computationally expensive the pooling layer, the better the accuracy. However, the use of the $G$-TC becomes prohibitive when $|G|$ increases. In the next section, we discuss the settings where each invariant layer should be preferred, and highlight each invariant layer's strengths and weaknesses.
\begin{table*}
\centering
\small
\begin{tabular}{||c|c|c|c|c|c|c|c||}
	\hline
	Dataset&Group&$G$&Pooling&$K$ filters&Avg acc.&Std. dev.&Param. count\\
 \hline
	\hline
	\multirow{8}{*}{MNIST}&\multirow{4}{*}{$\mathrm{SO}(2)$}&\multirow{4}{*}{$C_8$}&Avg $G$-pooling&24&0.74 &$<0.01$ &50247\\
	     &       &       &Max $G$-pooling&24& 0.96& $<0.01$ &50247\\
	     &       &       &Select. $G$-Bispectrum&24&0.95 &$<0.01$  &49116\\
	     &       &       &$G$-TC&24&0.96 &  $<0.01$&48385\\
        \cline{2-8}
	     &\multirow{4}{*}{$\mathrm{O}(2)$}&\multirow{4}{*}{$D_8$}&Avg $G$-pooling&4& 0.60&$<0.01$ &147675\\
	     &       &       &Max $G$-pooling&4&0.78 &  $<0.01$&147675\\
	     &        &      &Select. $G$-Bispectrum&4&0.93 & $<0.01$ &143029\\
	     &        &      &$G$-TC&4&0.96 &  $<0.01$&142220\\
	\hline
	\multirow{8}{*}{EMNIST}&\multirow{4}{*}{$\mathrm{SO}(2)$}&\multirow{4}{*}{$C_8$}&Avg $G$-pooling&24&0.40 &$<0.01$ &50195\\
	     &       &       &Max $G$-pooling&24& 0.76& $<0.01$ &50195\\
	     &       &       &Select. $G$-Bispectrum&24&0.77 & $<0.01$ &49254\\
	     &       &      &$G$-TC&24&0.80 &  $<0.01$&48494\\
	\cline{2-8}
	     &\multirow{4}{*}{$\mathrm{O}(2)$}&\multirow{4}{*}{$D_8$}&Avg $G$-pooling&20&0.38 &$<0.01$ &48832\\
	     &      &        &Max $G$-pooling&20&0.71 & $<0.01$ &48832\\
	     &      &        &Select. $G$-Bispectrum&20&0.74 & $<0.01$ &47320\\
	     &     &        &$G$-TC&20&0.79 & $<0.01$ &46954\\
	\hline
	%\multirow{8}{*}{CIFAR-10}&\multirow{4}{*}{$\mathrm{SO}(2)$}&Avg $G$-pooling&20&0.43 &$<0.01$ &10494140\\
	     %&              &Max $G$-pooling&20&0.48 & $<0.01$ &10494140\\
	     %&              &Select. $G$-Bsp&20& 0.50& $<0.01$ &10359300\\
	    % &              &$G$-TC&20&0.53 & $<0.01$ &10322860\\
	%\cline{2-7}
	     %&\multirow{4}{*}{$\mathrm{O}(2)$}&Avg $G$-pooling&20& 0.37&$<0.01$ &10494140\\
	    % &              &Max $G$-pooling&20& 0.43& $<0.01$ &10494140\\
	    % &              &Select. $G$-Bsp&20&0.38 &$<0.01$  &10316120\\
	    % &              &$G$-TC&20&0.48 & $<0.01$ &10276150\\
	     %\hline
\end{tabular}
\caption{Results of numerical experiments averaged over 10 runs with Avg $G$-pooling, Max $G$-pooling, our selective $G$-Bispectrum and $G$-TC. The experiments are performed on $\mathrm{SO}(2)/\mathrm{O}(2)$-MNIST and $\mathrm{SO}(2)/\mathrm{O}(2)$-EMNIST. % and -CIFAR-10.
The table shows the number of filters, the average classification accuracy, standard deviation and parameter count. This table shows that the selective $G$-Bispectrum conserves the accuracy of the $G$-TC at an equivalent number of parameters. \vspace{-0.1cm}}
\label{tab:classification}
\end{table*}
\paragraph{Discussion on the choice of invariant layer} The first observation from Table \ref{tab:classification} is though the \emph{selective} $G$-Bispectrum is complete, the model obtains slightly lower accuracy than $G$-TC. This observation might be surprising at first, since we prove mathematically in Section~\ref{sec:completeness} that the selective $G$-Bispectrum is complete, just as the full version. An explanation to this lies in the paradoxes of the Universal Approximation Theorem \citep{HORNIK1989359}. Just because an arbitrarily large MLP can theoretically fit any function, this does not imply that it will happen for a practical, limited MLP. In practice, we hypothesize that the redundancy of the $G$-TC allows the MLP to distinguish inputs more easily. If the size of the model allows it, the $G$-TC or the full $G$-Bispectrum will provide better accuracy. However, when the size of the group is big, their use is often out of reach while the selective $G$-Bispectrum is scalable.
In Table \ref{tab:classification}, we also notice that the Max $G$-pooling performs well compared to the others even though it is not complete. This is because we have many filters that allow for refined classification. Indeed, assume $f,\phi_k$ are black-and-white images with $N$ pixels. In consequence, $\max_g\Theta_k(g)\in\{0,1,...,N\}$ for $k=1,2,...,K$. The Max $G$-pooling allows a maximum separation of $(N+1)^K$ classes. In practice, this value is not reached, but it explains why Max $G$-pooling performs well. Figure \ref{fig:Kfilters} highlights this dependency of the Max $G$-pooling on the number of filters since the accuracy drops to less than $60\%$ with 2 filters. In comparison, the $G$-TC and the selective $G$-Bispectrum, which are \emph{complete}, keep an accuracy above $85\%$ with 2 filters.% If the parameter budget restricts the number of filters that can be implemented within the $G$-Conv, then we recommend using a complete invariant layer.

\begin{figure}
    \centering
    \includegraphics[width = 6.8cm]{Images/filters_MNIST_cyclic.pdf}
    \includegraphics[width = 6.8cm]{Images/filters_MNIST_dihedral.pdf}
    \includegraphics[width = 6.8 cm]{Images/filters_MNIST_cyclic2.pdf}
    \includegraphics[width = 6.8 cm]{Images/filters_MNIST_dihedral2.pdf}
    \vspace{-0.2cm}
    \caption{At the top: Evolution of the average classification accuracy with rotated MNIST ($\mathrm{SO}(2)$-MNIST) and rotated-reflected MNIST ($\mathrm{O}(2)$-MNIST) over $10$ runs when the number of filters varies from $2$ to $20$ for the Avg $G$-pooling, the Max $G$-pooling, the selective $G$-Bispectrum and the $G$-TC. The number of parameters of each model is maintained equal for fair comparison. The standard deviations are represented using vertical intervals. With the selective $G$-Bispectrum layer, we can reduce the number of convolutional filters needed for a given accuracy. For example, with only \( K = 2 \) filters, we achieve 96\% accuracy, compared to 63\% with the Max $G$-pooling layer. Our approach allows $G$-CNNs to maintain competitive accuracy while using smaller neural networks.  At the bottom, the same results are displayed with time instead of the number of filters on the $x$-axis. The dotted lines reproduce the evolution of $K$ from the figures at the top. We can observe that the selective $G$-Bispectrum is faster than the $G$-TC when a FFT is available, thus here in the case of $\mathrm{SO}(2)$-MNIST. Recall that an FFT can be implemented for many groups \cite{Diaconis1990EfficientCO}}\vspace{-0.5cm}
    \label{fig:Kfilters}
\end{figure}
\paragraph{Completeness} To conclude our numerical experiments, we study the robustness of the selective $G$-Bispectrum to adversarial attacks, following the analysis in \citet[Figure~2]{sanborn2023general}. Given an image $\widehat{f}:X\mapsto\mathbb{R}^c$ and a filter $\phi:X\mapsto\mathbb{R}^c$, they numerically verified  the robustness (=completeness) of the $G$-TC by showing that
\begin{equation}\label{eq:opti_sanborn}
    f^*\in\mathrm{arg}\min_{f:X\mapsto\mathbb{R}^c} \|T(\phi * f)-T(\phi * \widehat{f})\|_2^2 \iff f^*=\alpha(g, \widehat{f}) \text{ for some } g\in G.
\end{equation}
Indeed, \citet[Figure~2]{sanborn2023general} shows that only images that are identical up to rotation/reflection can yield the same $C_n/D_n$-TC. That is, the $G$-CNN with $G$-TC can not be ``fooled'' since only input in the same orbit yield the same output. We perform a similar experiment in Figure~\ref{fig:recovering_MNIST}. Moreover, it is well-known that the $G$-convolution $\phi*f$ is $G$-equivariant. Hence, an equivalent experiment is to show that
\begin{equation*}
    \Theta^*\in\mathrm{arg}\min_{\Theta:G\mapsto\mathbb{R}}\|T(\Theta)-T(\widehat{\Theta})\|^2_2\iff \Theta^* = \alpha(g, \widehat{\Theta})\text{ for some } g\in G.
\end{equation*}
In Figure~\ref{fig:completeness}, we show that the selective $G$-Bispectrum $\beta_{sel}$ is robust to adversarial attacks by solving
\begin{equation}   \label{eq:optimisation_prob}
    \Theta^*\in\mathrm{arg}\min_{\Theta:G\mapsto\mathbb{R}}\|\beta_{sel}(\Theta)-\beta_{sel}(\widehat{\Theta})\|^2_2.
\end{equation}
The signals are indeed recovered up to a translation, i.e., a group action of $C_{30}$. Moreover, despite~\eqref{eq:optimisation_prob} only optimizes using the selective $G$-Bispectrum, the full $G$-Bispectrum is correctly recovered. This is an additional evidence of the completeness of the selective $G$-Bispectrum.

\begin{figure}
\centering
\includegraphics[width=0.99\textwidth]{Images/Recovering_MNIST_V3bis.pdf}
\vspace{-0.3cm}
\caption{Adversarial attacks experiments with $G=C_4$. Images are optimized to output, respectively, a target selective $G$-bispectrum and a target $G$-Max pooling, obtained from an original image. The initial image is the initialization of the optimization process \eqref{eq:opti_sanborn}. After training, only for the selective $G$-Bispectrum (in blue), the recovered image is a copy of the original image up to group action (rotation). This is a numerical illustration of the robustness of the selective $G$-Bispectrum to adversarial attacks: one can not obtain the same output with an input that is not in the same class. On the other hand, $G$-Max pooling (in red) outputs a noisy image because it is not complete.}
\label{fig:recovering_MNIST}
\end{figure}

\begin{figure}[ht]
    %\includegraphics[width = 6cm]{Images/Recovered_signal.pdf}
    %\includegraphics[width = 6cm]{Images/Recovered_Bispectrum.pdf}
    \centering
    \includegraphics[width = 0.99\textwidth]{Images/Recovered_Bispectrum_and_signal_V2.pdf}
    \vspace{-0.2cm}
    \caption{Numerical experiment of signal recovering from original signals $\{\widehat{\Theta}_i\}_{i=1}^{15}$ where $\widehat{\Theta}_i:G\mapsto\mathbb{R}$ by solving \eqref{eq:optimisation_prob} with  $G = C_{30}$. We use the gradient method with Armijo line search to solve \eqref{eq:optimisation_prob}. The recovered solutions $\{\Theta^*_i\}_{i=1}^{15}$ are represented and correspond to translations of the original signals. The moduli of the full $G$-Bispectra are also represented and are identical. This experiment corroborates the completeness of the selective $G$-Bispectrum since we are able to recover an unknown signal only from the knowledge of its selective $G$-Bispectrum.
    }
    \label{fig:completeness}
\end{figure}

\section{Conclusion and Future works}
In this paper, we introduced a new type of complete invariant layer for $G$-invariant CNNs -- called \emph{selective $G$-Bispectrum layer} -- with the objective of increasing the accuracy and robustness of $G$-CNNs compared to those implemented with the initially proposed Max $G$-pooling. The $G$-TC layer also achieves this goal, but at an output cost of $\mathcal{O}(|G|^2)$ coefficients and $\mathcal{O}(|G|^3)$ flops that prevents its application to large groups, while the selective $G$-Bispectrum layer only outputs $\mathcal{O}(|G|)$ coefficients. Building on the result of \citet{kakarala_thesis} for cyclic groups, we have shown that the completeness of the \textit{selective} $G$-Bispectrum layer holds for all commutative groups, all dihedral groups, the octahedral and full octahedral groups. In a suite of experiments, we provided a global picture of the strength and weaknesses of each invariant layer. We studied the performance in terms of training speed, classification accuracy and robustness to adversarial attacks. As a result, the selective $G$-Bispectrum paves the way for the development of complete invariant pooling layers that can accommodate larger group sizes and, hence, a larger set of symmetries.
%For other groups of interest, such as the chiral octahedral and full octahedral groups in the context of 3D image processing \citep{Lenz2009OctahedralTF}, we provide a good heuristic providing a sequence of coefficients that allow the sequential inversion of the Bispectrum up to a degree of freedom. It remains to show that the degree of freedom corresponds to the inherent indeterminacy. The method suggests that the $G$-Bispectrum is complete with $\mathcal{O}(|G|)$ coefficients on many other groups, including the two previously cited.
%Given the wide variety of non-commutative groups, a general analysis would be an important breakthrough.
%\paragraph{Reproducibility} The code will be made available upon publication.%The code is  available on the github repository \url{https://github.com/geometric-intelligence/g-invariance}.
%\paragraph{Impact statement} This research contributes to advancements in the field of Convolutional Neural Networks. In image processing, this tool may be used for purposes such as facial recognition and automated security.
% \section{Acknowledgments}
% This work was supported by the Fonds National de la Recherche Scientifique -- FNRS.
\newpage
\bibliographystyle{icml2024}
\bibliography{bispectrumbib.bib}

\newpage
\appendix

\section{Background on groups}\label{app:group_theory}
We introduce the fundamentals of group theory, which provide the foundation for the theory of $G$-CNNs. These notions can be found in \citep{steinberg2011representation}.
% A natural starting point is to define a group.
\begin{definition}\label{def:group}
    A \emph{group} is a pair $(G,\ \cdot)$ where $G$ is a set and $\cdot:G\times G\mapsto G$ is an associative multiplication such that there is an identity element $e\in G$ (i.e., for all $g\in G$, $e\cdot g=g\cdot e=e$) and, for all $g\in G$, there is an inverse $g^{-1}\in G$ such that $g^{-1}\cdot g = g\cdot g^{-1}=e$.
\end{definition}
A group is thus a set $G$ combined with a product $\cdot$ preserving the characteristics of $G$. Here, the established term ``product'' can be misleading. It denotes any operation which makes Definition~\ref{def:group} true given the set $G$. For instance, $(\mathbb{R},+)$ is a group. Another example is $\mathrm{GL}(\mathbb{R}^n)$, the $n\times n$ real invertible matrices, associated to the usual matrix product. This group is said to be \emph{non-commutative} since $A\cdot B\neq B\cdot A$ in general for $A,B\in\mathrm{GL}(\mathbb{R}^n)$. An important group for us is the set $\{0,1,...,n-1\}$ associated with addition modulo~$n$. It is usually written $\mathbb{Z}/n\mathbb{Z}$ and called the \emph{cyclic group} $C_n$. A single group can arise in different contexts under seemingly distinct forms. For instance, $\mathbb{Z}/4\mathbb{Z}$ and the rotations leaving the square unchanged in $\mathbb{R}^2$ are fundamentally the same object. This observation gives rise to \emph{representation theory}, a branch of group theory studying how the same abstract idea of a group can emerge under different forms.
\begin{definition}\label{def:representation}
    A \emph{representation} of a group $(G,\ \cdot)$ is a pair $(\rho, V)$ where $V$ is a vector space and $\rho:G\mapsto \mathrm{GL}(V)$ is a group homomorphism, i.e., for all $g,h\in G$, $\rho(g\cdot h) = \rho(g)\rho(h)$. If $V$ is equipped with an inner product and if for all $g\in G$ and all $u,v\in V$,  $\langle\rho(g)v,\rho(g)w\rangle =\langle u, v\rangle$, $\rho$ is \emph{unitary}.
\end{definition}
\begin{remark}
    Throughout this paper, we use the shorthand $G$ to refer to the group $(G,\cdot )$ and $\rho$ to refer to a representation $(\rho,\ V)$.
\end{remark}
To illustrate Definition~\ref{def:representation}, a representation of $C_n$ is given by the complex roots of unity, $\rho(k) = \exp\left(\frac{2\pi i }{n}k\right)$, on the complex one-dimensional vector space $V= \mathbb{C}$. Every group also admit the \emph{trivial} representation: $\rho_0(g)=1$ for all $g \in G$. There is a specific subset of these representations called the irreducible representations, \emph{irreps} for short, being those that can not be expressed in a more compact form. The irreps are fundamental objects of group theory since they allow us to define an invertible Fourier transform on finite groups. The irreps are therefore needed to define the $G$-Bispectrum -- i.e., the Fourier transform of the $G$-TC. The notion of irreps is derived from that of a $G$-invariant subspace, which we recall in Definition~\ref{def:invariantspace}.
\begin{definition}\label{def:invariantspace}
    Given a representation $(\rho, V)$, a subspace $W\subseteq V$ is $G$-invariant if $\rho(g)w\in W$ for all $g\in G,\ w\in W$.
\end{definition}
The formal definition of the irreps is then stated as the representations with no non-trivial invariant subspace.
\begin{definition}
    A non-zero representation $(\rho,V)$ of a group $G$ is irreducible if the only $G$-invariant subspaces of $V$ are $\{0\}$ and $V$ itself.
\end{definition}
A single group acting on different spaces will have different representations. However, one can reveal the similarity between these representations by the mean of an equivalence relation.
\begin{definition}\label{def:equivalence}
    Two representations $(\rho,V)$ and $(\varphi, W)$ are \emph{equivalent} if there exists an isomorphism $T:V\mapsto W$ such that for all $g\in G$, $\rho(g) T = T \varphi(g)$.
\end{definition}
For the interested reader, the invertibility property of the Fourier transform is a consequence of the concepts of Pontryagin duality (commutative groups) and Tannaka-Krein duality (non-commutative groups); see, e.g., \citep{joyal_ross_tannaka}.
Our proofs will also rely on the notion of generating set of $G$, which we introduce here.
\begin{definition}
    A \emph{generating set} $S$ of a group $(G, \cdot)$ is a subset $S\subset G$ such that every $g\in G$ can be expressed as a finite combination of the elements in $S$ and their inverses under the group action $\cdot$.
\end{definition}
\begin{remark}
    It can be shown that every group $G$ of size $|G|$ has a generating set of size at most $\log_2 |G|$.
\end{remark}

\paragraph{Group Actions} A group $(G,\ \cdot)$ represents a set of transformations such as rotations that can act on data such as images.% , see Appendix \ref{app:group_theory} for the mathematical definition.
We define formally how groups can indeed transform datasets through the concept of group action.
%Now that a group $G$ and its irreps are properly defined, it is important to define how $G$ acts on a given vector space~$X$.
% Take for instance an infinite black-and-white picture modeled by some function $f:\mathbb{R}^2\mapsto [0,1]$. The \emph{group action} may define how the image is flipped or rotated. We define the notion of group action more formally below.
\begin{definition}\label{def:groupaction}
    Given a group $(G,\ \cdot)$, a \emph{group action} $\alpha:G\times X\mapsto X$ is a function satisfying i) Identity: $\alpha(e, x)=x$, ii) Compatibility: $\alpha(h,\alpha(g,x))=\alpha(h\cdot g, x)$ for any $x \in X$ and $h, g \in G$ and where $e$ is the identity of $G$.
\end{definition}

%Group representations are particular types of group actions that act on vector spaces $V$: we denote them by $(\rho,\ V)$. Irreducible representation (irreps for short) are special group representations where the vector space $V$ has minimal dimension in a sense made precise in Appendix \ref{app:group_theory}.

Processing operations and neural networks can be designed so that they respect group actions: specifically, a group acting on the input (e.g., rotating an input image) should yield a group action on the output (e.g., a rotation of the output feature map). This is the notion of $G$-\emph{equivariance}.
\begin{definition}
    A function $\psi:X\mapsto Y$ is $G$-\emph{equivariant} if $\psi(\alpha_1(g,x)) = \alpha_2(g,\psi(x))$ for all $x\in X$ and all $g\in G$, where $\alpha_1$ and $\alpha_2$ are group actions on $X$ and $Y$, respectively.
\end{definition}
For example, the $G$-convolution layer is $G$-equivariant by design  \citep{cohenc16}. An important problem in signal processing and deep learning is to achieve \textit{invariance} to nuisance factors not relevant for the task. Many of these factors are describable as group actions (e.g. rotations, translations, scaling). Thus, we want processing methods and machine learning models to be $G$-invariant:
\begin{definition}
    A function $\psi:X\mapsto Y$ is $G$-\emph{invariant} if $\psi(\alpha(g,x)) = \psi(x)$ for all $x\in X$ and all $g\in G$.
\end{definition}
For example, the Max $G$-pooling ($\max_{g\in G}\Theta(g)$) traditionally follows a $G$-convolutional layer to remove the equivariance of the convolution and achieve $G$-invariance. A $G$-CNN is a neural network that consists of $G$-convolutional layers and a pooling/invariance operation. The main applications of our proposed selective $G$-Bispectrum operation is to act as a $G$-invariant pooling layer, that can conveniently replace the classical Max $G$-Pooling layer of $G$-CNNs, as shown in the rest of the paper.
\begin{figure}
    \centering
    \includegraphics[width = 7cm]{Images/completeness_figure.png}
    \caption{Illustration of the concepts of excessive and complete invariance to a group action. With the excessive invariance, samples from different classes can be mapped to the same output.}
    \label{fig:sophias_figure}
\end{figure}

\paragraph{Clebsh-Jordan matrices}Given a group $(G,\cdot)$ and a family of unitary irreps $\{\rho_i\}_{i=0}^{|\text{Irreps}|-1}$, the Clebsh-Jordan matrix is analytically defined for each pair $\rho_1, \rho_2$ as:
\begin{equation}
    (\rho_1 \otimes \rho_2)(g) = C_{\rho_1,\rho_2} \Big[ \bigoplus_{\rho \in\mathcal{R}} \rho(g) \Big] C_{\rho_1,\rho_2}^\dagger.
    \label{eqn:non-commutative-bs}
\end{equation}

\paragraph{$G$-Bispectrum for Commutative Finite Groups}

The computation of the $G$-Bispectrum simplifies for commutative groups compared to Theorem~\ref{def:Bispectrum}, as recalled below.

\begin{theorem}{\citep{kakarala_finite_groups}}\label{thm:commutative_Bispectrum}
    If $(G,\ \cdot)$ is a commutative group and $\Theta:G\rightarrow \mathbb{R}$, the $G$-Bispectrum can be computed as
    \begin{equation}\label{eq:com_bsp}
        \beta(\Theta)_{\rho_1,\rho_2} = \mathcal{F}(\Theta)_{\rho_1}\mathcal{F}(\Theta)_{\rho_2}\mathcal{F}(\Theta)_{\rho_1\otimes\rho_2}^\dagger.
    \end{equation}
\end{theorem}
For commutative groups, the $G$-bispectral coefficients are complex scalars \citep{steinberg2011representation}.

\section{Indeterminacy of $G$-Bispectrum inversion problem}\label{sec:indeterminacy}
It is important to state precisely which information we can possibly retrieve from the $G$-Bispectrum.
A consequence of the $G$-invariance of $\beta(\Theta)$ is that the $G$-Bispectrum inversion problem is ill-posed. Recall that $G$-invariance means that for all $h\in G$, we have $\beta(\alpha(h,\Theta))_{\rho_1,\rho_2} =\beta(\Theta)_{\rho_1,\rho_2}$. Given a function $\Theta:G\mapsto\mathbb{R}$, a possible definition for the group action on $\Theta$ is given by $\alpha(h, \Theta(g)) =\Theta(h^{-1}\cdot g)$ for all $h\in G$ (see, e.g., \citep{cohenc16}). Therefore, for all $h\in G$, we have
\begin{align*}
        \mathcal{F}(\alpha(h,\Theta))_\rho &= \sum_{g\in G} \alpha(h, \Theta(g)) \rho(g)^\dagger\\
        &=\rho(hh^{-1})\sum_{g\in G} \Theta(h^{-1}g) \rho(g)^\dagger\\
        &= \rho(h) \mathcal{F}(\Theta)_\rho,
    \end{align*}
which shows that the $G$-Fourier transform $\mathcal{F}(\Theta)$ is $G$-equivariant. In consequence, recovering $\mathcal{F}(\Theta)$ from $\beta(\Theta)$ can at best be done up to an unknown factor $\rho(h)$. Moreover, as explained in \citep{kakarala_thesis, kondor08}, the indeterminacy is not limited to $\rho(h)$. Take for instance $C_n$. An indeterminacy factor $\rho_k(h) = \exp\left(\frac{2\pi i h k}{n}\right)$ corresponds to a translation of $h\in C_n$ of the signal, $\widetilde{\Theta}(g)=\Theta(g+h)$. \citep{kakarala_thesis} showed that $h$ is not restricted to $C_n=\mathbb{Z}/n\mathbb{Z}$: it may take any value in $[0,n]$. The Bispectrum is not only invariant to a discrete set of rotations, but to the continuous group of rotations $\mathrm{SO}(2)$. The factor can thus be written $\exp(i \varphi k)$ where $\varphi\in[0, 2\pi)$.

\section{Selective $G$-Bispectrum inversion: known results}\label{sec:commutative_alg}
%We present the algorithm from \citep{kakarala_thesis} to invert the Bispectrum the cyclic group $C_n=$($\mathbb{Z}/n\mathbb{Z}$, $+$ mod $n$). The method requires only $|G|=n$ bispectral coefficients, therefore allowing the reduction in complexity expected with a bispectral layer compared to the TC layer. We prove mathematically that the method can be extended to all commutative groups, still requiring only $|G|$ bispectral coefficients. Moreover, we show that it can be extended to the dihedral group $D_n$, i.e., the symmetries of the $n$-gon. In this case, the method only requires $\frac{|G|}{4}$ bispectral coefficients.
From now on, we assume that the Fourier transform $\mathcal{F}(\Theta)$ only features \textbf{non-zero elements}, or \textbf{invertible matrices} in the case of non-scalar Fourier transform. This assumption is supported by the zero probability of encountering this corner case.

\paragraph{Cyclic groups $C_n$} We start with $(G,\ \cdot)=(\mathbb{Z}/n\mathbb{Z},\ +\ \mathrm{mod}\ n)=:C_n$. Recall that the irreps are given by $\rho_k(g)=\exp\left(\omega_k g\right)$ where $\omega_k:=\frac{2\pi i k }{n}$ for $k\in\mathbb{Z}/n\mathbb{Z}$ (see, e.g., \citep{steinberg2011representation}).
\begin{theorem}{\citep{kakarala_finite_groups}}\label{thm:cyclic_inversion}
For cyclic groups $C_n$, the $C_n$-Bispectrum can be inverted using $|G|=n$ coefficients if $\mathcal{F}(\Theta)_\rho\neq 0$ for all $\rho \in C_n$.
\end{theorem}
\begin{proof}
The Fourier coefficient associated to the trivial representation $\mathcal{F}(\Theta)_{\rho_0}$, is uniquely determined and can be recovered from Theorem~\ref{thm:commutative_Bispectrum} by identifying phase and modulus:
\begin{equation}\label{eq:trivial_coef}
    \beta(\Theta)_{\rho_0,\rho_0} =|\mathcal{F}(\Theta)_{\rho_0}|^3\exp\left(i\ \mathrm{arg}(\mathcal{F}(\Theta)_{\rho_0})\right).
\end{equation}

We proceed using Pontryagin duality: the irreps $\{\rho_k\}_{k=1}^n$, form a group $\widehat{G}$ themselves, with $\widehat{G} = G$. In the case of the cyclic group $C_n$, notice that for all $j,k\in\mathbb{Z}/n\mathbb{Z}$, $\rho_j\otimes \rho_k = \rho_{j+k}$. Leveraging \eqref{eq:com_bsp}, we can use $\beta(\Theta)_{\rho_0,\rho_1}$ to recover $\mathcal{F}(\Theta)_{\rho_1}$:

\begin{equation}\label{eq:trivial_coeff_cyclic}
    |\mathcal{F}(\Theta)_{\rho_1}|^2 =\frac{\beta(\Theta)_{\rho_0,\rho_1}}{\mathcal{F}(\Theta)_{\rho_0}}.
\end{equation}

Equation~\eqref{eq:trivial_coeff_cyclic} leaves an indeterminacy on the phase of $\mathcal{F}(\Theta)_{\rho_1}$. This corresponds to the indeterminacy factor $\exp(i \varphi)$, $\phi\in[0,2\pi)$ of Appendix~\ref{sec:indeterminacy}. It is inherited from the $G$-invariance  of $\beta(\Theta)$ (it is not injective, hence you cannot distinguish inputs that have the same $G$-Bispectrum). For now, let $\arg(\mathcal{F}(\Theta)_{\rho_1})=0$. The key to recover all the other Fourier coefficients is to notice that $S = \{1\}$ is a generating set of $\widehat{G} =\mathbb{Z}/n\mathbb{Z}$. Therefore, computing sequentially
\begin{equation}
    \mathcal{F}(\Theta)_{\rho_{k+1}} = \left(\frac{\beta(\Theta)_{\rho_1,\rho_k}}{\mathcal{F}(\Theta)_{\rho_1}\mathcal{F}(\Theta)_{\rho_k}}\right)^\dagger,
\end{equation}
for $k=1,2,...,n-2$ recovers completely $\mathcal{F}(\Theta)$. We are not done yet because the phase we fixed before is not a valid shift. A valid phase shift for $\mathcal{F}(\Theta)_{\rho_1}$ is such that the shift w.r.t the original signal has the form $\exp\left(\frac{2\pi i h}{n}\right)$ for $h\in\mathrm{N}$. This valid phase shift is easy to find. It is the unique $\varphi\in[0,\frac{2\pi}{n})$ such that, if we define $\mathcal{F}(\widetilde{\Theta})_{\rho_k} = \exp(\varphi k) \mathcal{F}(\Theta)_{\rho_{k}}$ for all $k\in\mathbb{Z}/n\mathbb{Z}$, then we have $\mathcal{F}^{-1}(\mathcal{F}(\widetilde{\Theta}))\in\mathbb{\mathbb{R}}^n$ (i.e., with no imaginary part). Note that this is an explicit method to recover a valid phase while \citep{kakarala_finite_groups} relies on its existence without explicit method to find it.
The method is summarized in Algorithm~\ref{alg:cyclic_alg} and illustrated in Figure \ref{fig:alg1}. In consequence only the following $G$-bispectral coefficients are needed for completeness: $\beta(\Theta)_{\rho_0,\rho_0}, \beta(\Theta)_{\rho_0,\rho_1}$ and $ \beta(\Theta)_{\rho_1,\rho_k}$ for $k=1,2,...,n-2$. This makes a total of $n=|G|$ coefficients. We summarize this result in Theorem~\ref{thm:cyclic_inversion}.
\end{proof}


\begin{figure}
\begin{minipage}[c]{0.54\textwidth}
  \begin{algorithm}[H]
    \caption{Bispectrum inversion on $\mathbb{Z}/n\mathbb{Z}$
 \citep{kakarala_finite_groups}}
    \label{alg:cyclic_alg}
    \small
    \begin{algorithmic}[1]
        \STATE \textbf{Input}: $\beta(\Theta)_{\rho_0,\rho_0}, \beta(\Theta)_{\rho_0,\rho_1}$ and $ \beta(\Theta)_{\rho_1,\rho_k}$ for $k=1,2,...,n-2$.
        \STATE Compute $|\mathcal{F}(\Theta)_{\rho_0}|=\left(\beta(\Theta)_{\rho_0,\rho_0}\right)^{\frac{1}{3}}$ and $\mathrm{arg}(\mathcal{F}(\Theta)_{\rho_0})=\mathrm{arg}(\beta(\Theta)_{\rho_0,\rho_0})$.
        \STATE Compute $|\mathcal{F}(\Theta)_{\rho_1}| =\left(\frac{\beta(\Theta)_{\rho_0,\rho_1}}{\mathcal{F}(\Theta)_{\rho_0}}\right)^{\frac{1}{2}}$ and set $\mathrm{arg}(\mathcal{F}(\Theta)_{\rho_1})=0$.
        \FOR{$k=1,2,...,n-2$}
            \STATE Compute $\mathcal{F}(\Theta)_{\rho_{k+1}} = \left(\frac{\beta(\Theta)_{\rho_1,\rho_k}}{\mathcal{F}(\Theta)_{\rho_1}\mathcal{F}(\Theta)_{\rho_k}}\right)^\dagger$.
        \ENDFOR
        \STATE Find $\varphi\in [0,\frac{2\pi}{n})$ such that $\mathcal{F}^{-1}(\mathcal{F}(\widetilde{\Theta}))\in\mathbb{R}^n$ where $\mathcal{F}(\widetilde{\Theta})_{\rho_k} = \exp(\phi k)\mathcal{F}(\Theta)_{\rho_{k}}$.
        \STATE \textbf{Return} $\mathcal{F}(\widetilde{\Theta})$ (up to group action).
    \end{algorithmic}
\end{algorithm}
\end{minipage}\hfill%
\begin{minipage}[c]{0.4\textwidth}
    \includegraphics[width = 6cm]{Images/alg1_icml.pdf}
    \caption{Illustration of Algorithm \ref{alg:cyclic_alg}. The Bispectrum coefficients allow to recover the Fourier transform sequentially, up to group action. First, $\beta_{\rho_0,\rho_0}$ gives $\mathcal{F}_{\rho_{0}}$, which, combined with $\beta_{\rho_0,\rho_1}$ gives $\mathcal{F}_{\rho_{1}}$ etc. }
    \label{fig:alg1}
\end{minipage}
\end{figure}

\section{Selective $G$-Bispectrum inversion: commutative groups}\label{app:commutative_inversion}

Here, we prove the theorem stated in the main text and recalled below:
\begin{theorem}
For finite commutative groups $G$, the $G$-Bispectrum can be inverted using $|G|$ coefficients if $\mathcal{F}(\Theta)_\rho \neq 0$ for all $\rho \in G$.
\end{theorem}


Specifically, we extend Algorithm~\ref{alg:cyclic_alg} to all commutative groups, based on the Theorem~\ref{thm:fundamentalthm}. That is, we design a method for the direct sum of finitely many cyclic groups.
\begin{theorem}{(see, e.g., \citep{Norman_2012})}\label{thm:fundamentalthm}
    Every finite commutative group $G$ is isomorphic to a finite direct sum of cyclic groups: $G\cong \bigoplus_{l=1}^L \mathbb{Z}/n_l\mathbb{Z}$ where $L\in\mathbb{N}$ and $n_l\in\mathbb{N}$ for $l=1,2,...,L$.
\end{theorem}
For all $\kbf\in G$ ($\kbf$ is integer-valued vector of length $L$), the irreps $\rho_\kbf:G\mapsto \mathbb{C}$ are given by $
    \rho_\kbf(g) = \prod_{l=1}^L\exp(\frac{2\pi i\kbf_l}{n_l} g_l)$.
The number of irreps is $|G| = \prod_{l=1}^L n_l$. We detail and prove in Theorem~\ref{thm:invert_commutative} the procedure to invert the $G$-Bispectrum on commutative groups. The procedure is summarized in Algorithm~\ref{alg:commutative_alg} where we use the two following notations.\begin{enumerate}
    \item $\ebf^l$ denotes the basis vector in $\mathbb{Z}^L$ such that $\ebf^l_k=1$ if $k=l$ and $\ebf^l_k=0$ otherwise.
    \item $K^l:=\{t\ebf^l +\kbf \ | \ t=1,2...,n_l-1\text{ and } \kbf\in K^{l-1}$\} for $l=1,2,...,L$.
\end{enumerate}
The sets $K^l$ are a recursively constructed such that $K^L\cong G$. For $G = (\mathbb{Z}/3\mathbb{Z})^3$, the sets $K_1,K_2,K_3$ are represented in Figure~\ref{fig:ksets}.
\begin{figure}[h]
    \centering
    \includegraphics[width = 5cm]{Images/K_sets.pdf}
    \caption{Representation of the sets $K_1,K_2$ and $K_3$ for $G = (\mathbb{Z}/3\mathbb{Z})^3$.}
    \label{fig:ksets}
\end{figure}
\begin{theorem}\label{thm:invert_commutative}
For finite commutative groups $G$, the $G$-Bispectrum can be inverted using $|G|$ coefficients if $\mathcal{F}(\Theta)_\rho \neq 0$ for all $\rho \in G$.
\end{theorem}
\begin{proof}
%For all $\kbf\in G$, the irreps $\rho_\kbf:G\mapsto \mathbb{C}$ are given by $\rho_\kbf(g) = \prod_{l=1}^L\exp(\frac{2\pi i \kbf_l}{n_l} g_l)$. The number of irreps is $|G| = \prod_{l=1}^L n_l$.
Notice that for the commutative groups, we keep the property $\rho_{\kbf}\otimes\rho_{\kbf'}=\rho_{\kbf+\kbf'}$ where $(\kbf+\kbf')_l := \kbf_l+\kbf_l'\ \mathrm{mod} \ n_l$ for $l=1,2,...,L$. The first step is to obtain a generating set $S$ of the irreps is of size $L$. It is given by the usual basis vectors $S=\{\ebf^l\in \mathbb{Z}^L$ for $l=1,...,L\}$ where
\begin{equation*}
    \ebf^l_k=\begin{cases}
        1 \text{ if } k=l,\\
        0 \text{ otherwise.}
    \end{cases}
\end{equation*}
By Theorem~\ref{thm:commutative_Bispectrum}, it is sufficient to have the Fourier coefficients associated to each generating element in $S$ to recover all the Fourier coefficients. Indeed, knowing $\mathcal{F}(\Theta)_{\rho_{\kbf_1}}$ and $\mathcal{F}(\Theta)_{\rho_{\kbf_2}}$ allows us to compute $\mathcal{F}(\Theta)_{\rho_{\kbf_1+\kbf_2}}$. By definition of the generating set $S$, we can thus recover $\mathcal{F}(\Theta)_{\rho_{\kbf}}$ for all $\kbf \in G$. The moduli of the coefficients $\mathcal{F}(\Theta)_{\rho_{\ebf^l}}$ for $l=1,2,...,L$ can be computed as follows:
\begin{equation}\label{eq:compute_generating}
    |\mathcal{F}(\Theta)_{\rho_{\ebf^l}}| =\left(\frac{\beta(\Theta)_{\rho_\mathbf{0},\rho_{\ebf^l}}}{\mathcal{F}(\Theta)_{\rho_\mathbf{0}}}\right)^\frac{1}{2},
\end{equation}
where $\mathcal{F}(\Theta)_{\rho_\mathbf{0}}$ is the Fourier coefficient of the trivial representation (computed as in~\eqref{eq:trivial_coef}). We claim that the phase can be fixed independently for each label $l$, thus $L$ times. This is because only one factor $h_l$ remains among the $L$ independent factors in $\mathbf{h}$:
\begin{align*}
    \mathcal{F}(\alpha(\mathbf{h},\Theta))_{\rho_{\ebf^l}} &= \rho_{\ebf^l}(\mathbf{h})\mathcal{F}(\Theta)_{\rho_{\ebf^l}}=\exp\left(\frac{2\pi i h_l}{n_l}\right)\mathcal{F}(\Theta)_{\rho_{\ebf^l}},
\end{align*}
for all $l=1,2,.,L$. Therefore, fixing an arbitrary phase in~\eqref{eq:compute_generating} only fixes $h_l$. Again, the indeterminacy factor $h_l$ is not restricted to $\mathbb{Z}/n_l\mathbb{Z}$ but can belong to $[0,n_l]$. We will have to solve this issue further. For now, we set the phase of $\mathcal{F}(\Theta)_{\rho_{\ebf^l}}$ to zero. Once the Fourier coefficients are known for all the generators of the group of the irreps $\{\rho_{\ebf^l}\}_{l=1}^L$, it remains to combine them to obtain all the elements in the groups and, consequently, all the associated Fourier coefficients.

At this point, it helps to consider the problem geometrically. Each irreps $\rho_\kbf$ can be associated to its integer coordinate $\kbf$ inside a hyper-rectangle in $\mathbb{R}^L$, whose length of edges is $n_l$ for $l=1,2,...,L$. We combine the coordinates $\ebf^l$ to obtain all the possible integer coordinates inside the hyper-rectangle. First, we can obtain the $L$ orthogonal edges of the hyper-rectangle. For $l=1,2,...,L$, $\rho_{\ebf^l}\otimes \rho_{\ebf^l}$ gives $\rho_{2\ebf^l}$, $\rho_{\ebf^l}\otimes \rho_{2\ebf^l}$ gives $\rho_{3\ebf^l}$, etc. This is in fact the procedure of Algorithm 1. Now, we combine the edges to generate the inside of the hyper-rectangle. We proceed iteratively. For $l=1,2,...,L$, we define $K^l:=\{t\ebf^l + \kbf \ | \ t=1,2,..,n_l-1\text{ and } \kbf\in K^{l-1}\}$ and $K^0:=\{\mathbf{0}\}$. This construction is such that $K^l\cong\bigoplus_{j=1}^l \mathbb{Z}/n_j\mathbb{Z}$ and $K^L=G$. We generate the missing Fourier coefficients by combining the ones associated to the generating set of $G$.
For $l=2,...,L$, compute
\begin{equation}
    \mathcal{F}(\Theta)_{\rho_{t\ebf^l+\kbf}}=\left(\frac{\beta(\Theta)_{\rho_{t\ebf^l},\rho_{\kbf}}}{\mathcal{F}(\Theta)_{\rho_{t\ebf}}\mathcal{F}(\Theta)_{\rho_{\kbf}}}\right)^\dagger,
\end{equation}
for $t=1,2,...,n_l-1$, for all $\kbf\in K^{l-1}$. Intuitively for $L=3$, we obtain first an edge, then a face and finally the full parallelepiped. To conclude, we reproduce the procedure from Algorithm~\ref{alg:cyclic_alg} to find a valid phase shift $\varphi_l$ in each basis direction $\ebf^l$. The last step is then to compute, for all $\kbf\in G$, $$\mathcal{F}(\widetilde{\Theta})_{\rho_{\kbf}} =\mathcal{F}(\Theta)_{\rho_{\kbf}} \prod_{l=1}^L\exp(\phi_l k_l).$$  The procedure is summarized in Algorithm~\ref{alg:commutative_alg} and illustrated in Figure \ref{fig:alg2}. It shows that the bispectral coefficients needed for completeness are: $\beta(\Theta)_{\rho_0,\rho_0}$, $\beta(\Theta)_{\rho_0,t\rho_{\ebf}^l},\ \beta(\Theta)_{t\rho_{\ebf}^l, \rho_\kbf} $ for $l=1,2,...,L$, $t=1,2,...,n_l-2$ and all $\kbf\in K^{l-1}$. We recover exactly one Fourier coefficient per $G$-bispectrum coefficient. This makes thus a total of $|G|$ bispectral coefficients precisely and proves the following theorem.
\end{proof}

\begin{figure}
\begin{minipage}[c]{0.54\textwidth}
\begin{algorithm}[H]
    \caption{Bispectrum inversion for finite commutative groups ($G=\bigoplus_{l=1}^L \mathbb{Z}/n_l\mathbb{Z}$).}
    \label{alg:commutative_alg}
    \small
    \begin{algorithmic}[1]
        \STATE \textbf{Input}: $|G|$ bispectral coeffs.
        \STATE Compute $|\mathcal{F}(\Theta)_{\rho_\mathbf{0}}|=\left(\beta(\Theta)_{\rho_\mathbf{0},\rho_\mathbf{0}}\right)^{\frac{1}{3}}$ and $\mathrm{arg}(\mathcal{F}(\Theta)_{\rho_\mathbf{0}})=\mathrm{arg}(\beta(\Theta)_{\rho_\mathbf{0},\rho_\mathbf{0}})$.
        \FOR{$l=1,...,L$}
            \STATE Compute $|\mathcal{F}(\Theta)_{\rho_{\ebf^l}}| =\left(\frac{\beta(\Theta)_{\rho_\mathbf{0},\rho_{\ebf^l}}}{\mathcal{F}(\Theta)_{\rho_\mathbf{0}}}\right)^\frac{1}{2}$ and set $\mathrm{arg}\left(\mathcal{F}(\Theta)_{\rho_{\ebf^l}}\right)=0$.
            \FOR{$t=1,...,n_l-1$}
                \IF{$t>1$}
                \STATE  $F(\Theta)_{\rho_{t\ebf^l}}=\left(\frac{\beta(\Theta)_{\rho_{\ebf^l},\rho_{(t-1)\ebf^l}}}{\mathcal{F}(\Theta)_{\rho_{\ebf^l}}\mathcal{F}(\Theta)_{\rho_{(t-1)\ebf^l}}}\right)^\dagger$.
                \ENDIF
                \FOR{$\kbf\in K^{l-1}\setminus\{\mathbf{0}\}$}
                    \STATE Compute $\mathcal{F}(\Theta)_{\rho_{t\ebf^l+\kbf}}=\left(\frac{\beta(\Theta)_{\rho_{t\ebf^l},\rho_{\kbf}}}{\mathcal{F}(\Theta)_{\rho_{t\ebf^l}}\mathcal{F}(\Theta)_{\rho_{\kbf}}}\right)^\dagger$.
                \ENDFOR
            \ENDFOR
        \ENDFOR
        \FOR{$l=1,2,...,L$}
            \STATE Find $\varphi_l \in [0,\frac{2\pi}{n})$ s.t. $\mathcal{F}^{-1}(\mathcal{F}(\widetilde{\Theta})_{\rho_{(1,..,n_l)\ebf^l}})\in\mathbb{R}^{n_l}$ where $\mathcal{F}(\widetilde{\Theta})_{\rho_{k\ebf^l}} = \exp(\phi_l k)\mathcal{F}(\Theta)_{\rho_{k\ebf^l}}$.
        \ENDFOR
        \FOR{$\kbf\in G$}
           \STATE $\mathcal{F}(\widetilde{\Theta})_{\rho_{\kbf}} =\mathcal{F}(\Theta)_{\rho_{\kbf}} \prod_{l=1}^L\exp(\phi_l k_l)$.
        \ENDFOR
        \STATE \textbf{Return} $\mathcal{F}(\widetilde{\Theta})$ (up to group action).
    \end{algorithmic}
\end{algorithm}
\end{minipage}\hfill
\begin{minipage}[c]{0.4\textwidth}
    \centering
    \includegraphics[width = 6cm]{Images/alg2_icml.pdf}
    \caption{Illustration of Algorithm \ref{alg:commutative_alg}. The Bispectrum coefficients allow to recover the Fourier transform sequentially, up to group action.}
    \label{fig:alg2}
\end{minipage}
\end{figure}

\section{Selective $G$-Bispectrum inversion: dihedral groups}\label{app:dihedral_inversion}

\paragraph{Dihedral group $D_n$}\label{sec:inversion_dihedral}
The dihedral group $D_n$ is the group of all symmetries of the $n$-gon. Mathematically, it is defined as
\begin{equation}\label{def:dihedral}
    D_n:=\langle x,\ a\ |\ a^n=x^2=e,\ xax = a^{-1}\rangle,
\end{equation}
where $a$ is the \emph{rotation} and $x$ is the \emph{reflection}, and they form a generating set of $D_n$. We will only consider the case $n>2$ since the cases $n=1$ and $n=2$ are commutative groups covered by the previous subsection, while $n>2$ gives non-commutative groups. The 2D irreps of $D_n$ are given by
\begin{equation}\label{eq:2d_irreps_dn}
    \rho_k(a^lx^m)=\begin{bmatrix}
        \cos(\omega_l k)&-\sin(\omega_l k)\\
        \sin(\omega_l k)&\cos(\omega_l k)
    \end{bmatrix}
    \begin{bmatrix}
        1 & 0\\
        0 & -1\\
    \end{bmatrix}^m,
\end{equation}
where $\omega_l = \frac{2\pi l}{n}$ for $k=1,2,...,\lfloor \frac{n-1}{2}\rfloor$. There are also $2$ or $4$ 1D irreps if $n$ is odd or even, respectively. We denote these 1D irreps by $\rho_{0},\ \rho_{01},\ \rho_{02}$, and $\rho_{03}$ (see Appendix \ref{app:1dirreps}). The two last ones only exist for $n$ even.
\begin{theorem}\label{thm:dihedral_inversion}
    For the family of dihedral groups $D_n$, we need at most $\left \lfloor{ \frac{n-1}{2}}\right \rfloor + 2$ bispectral matrix coefficients for inversion if $\det(\mathcal{F}(\Theta)_\rho) \neq 0$ for all irreps $\rho$ of $D_n$. This corresponds to $1+4+16\cdot\left\lfloor{ \frac{n-1}{2}}\right \rfloor\approx 4 |D_n|$ scalar values.
\end{theorem}
\begin{proof}
In view of section~\ref{sec:commutative_alg}, we wish to show that there is an irrep $\rho_1$ that generates all the irreps of $D_n$. As in the cyclic and commutative cases, we can first deduce $\mathcal{F}(\Theta)_{\rho_0}$ from $\beta(\Theta)_{\rho_0,\rho_0}$. Now, we have
\begin{equation}\label{eq:first_coeff_dihedral}
    \mathcal{F}(\Theta)_{\rho_1}\mathcal{F}(\Theta)_{\rho_1}^\dagger=\frac{\beta(\Theta)_{\rho_0,\rho_1}}{\mathcal{F}(\Theta)_{\rho_0}}.
\end{equation}
The novelty for this non-commutative group is that $\mathcal{F}(\Theta)_{\rho_1}$ is a $2\times 2$ matrix. After computing the eigenvalue decomposition $\frac{\beta(\Theta)_{\rho_0,\rho_1}}{\mathcal{F}(\Theta)_{\rho_0}} = V\Lambda V^\dagger$, we can choose
\begin{equation}
    \mathcal{F}(\Theta)_{\rho_1} = V\Lambda^{\frac{1}{2}}V^\dagger U.
\end{equation}
for all unitary $U$ (i.e., $UU^\dagger=U^\dagger U=I$) to solve~\eqref{eq:first_coeff_dihedral}. In the case $\mathbb{Z}/n\mathbb{Z}$ the indeterminacy belonged to $\mathrm{SO}(2)$, the continuous set of 2d rotations. For $D_n$, it belongs to $\mathrm{O}(2)$, the continuous set of 2d rotations and reflections. For finite groups, \citet{kakarala_thesis} ensures that the only choices for $U$ such that $\mathcal{F}(\Theta)$ is a Fourier transform on $G(=D_n)$ are such that $\Theta$ is identical to the original signal up to some group action $g\in D_n$. %that respect $\eqref{eqn:non-commutative-bs}$ are such that $U = \rho_k(a^l x^m)$ for some $k,l,m$ (see \eqref{eq:2d_irreps_dn}).
For $k=2,...,\lfloor \frac{n-1}{2}\rfloor$, we can then obtain
\begin{align}
    \label{eq:new_coeff_dihedral}
    &\bigoplus_{\rho\in\rho_1\otimes\rho_{k-1}}\mathcal{F}(\Theta)_\rho =\left(C_{\rho_1,\rho_{k-1}}^{\dagger}\left[\mathcal{F}(\Theta)_{\rho_1}\otimes\mathcal{F}(\Theta)_{\rho_{k-1}}\right]^{-1}\beta(\Theta)_{\rho_1,\rho_2}C_{\rho_1,\rho_{k-1}}\right)^\dagger.
\end{align}
It is shown in Appendix~\ref{app:generating_dihedral} that $\rho_k$ appears in the tensor decomposition \eqref{eqn:non-commutative-bs} of $\rho_1\otimes\rho_{k-1}$ so that the for-loop can be applied. Moreover, we know from Appendix~\ref{app:generating_dihedral} that $\rho_{01}$ appears in the decomposition of $\rho_1\otimes\rho_{1}$ and, for $n$ even, $\rho_{02}, \rho_{03}$ in $ \rho_1\otimes\rho_{\frac{n}{2}-1}$. Thus the iteration recovers the complete DFT $\mathcal{F}(\Theta)$. $\beta(\Theta)_{\rho_0,\rho_0}$ is a scalar, $\beta(\Theta)_{\rho_0,\rho_1}$ is a $2\times 2$ matrix and $ \beta(\Theta)_{\rho_1,\rho_k}$, $k\neq 0$, is a $4\times 4$ matrix. Hence the total number of scalars that is required is $1+4+16\lfloor\frac{n-1}{2}\rfloor$.
\end{proof}
The procedure is summarized in Algorithm \ref{alg:dihedral_alg}.
\begin{algorithm}[tb]
    \caption{Bispectrum inversion on the dihedral group $D_n$.}\label{alg:dihedral_alg}
    \small
    \begin{algorithmic}[1]
        \STATE \textbf{Input}: $\beta(\Theta)_{\rho_0,\rho_0}, \beta(\Theta)_{\rho_0,\rho_1}$ and $ \beta(\Theta)_{\rho_1,\rho_k}$ for $k=1,2,...,\lfloor\frac{n-1}{2}\rfloor$.
        \STATE Compute $|\mathcal{F}(\Theta)_{\rho_0}|=\left(\beta(\Theta)_{\rho_0,\rho_0}\right)^{\frac{1}{3}}$ and $\mathrm{arg}(\mathcal{F}(\Theta)_{\rho_0})=\mathrm{arg}(\beta(\Theta)_{\rho_0,\rho_0})$.
        \STATE Compute $V\Lambda V^\dagger = \frac{\beta(\Theta)_{\rho_0,\rho_1}} {\mathcal{F}(\Theta)_{\rho_0}}$.\\
        \STATE Set $\mathcal{F}(\Theta)_{\rho_1} = V\Lambda^\frac{1}{2} V^\dagger U$ with valid $U$ (existence of $U$ ensured but not computed easily).\\
        \FOR{$k=2,...,\lfloor\frac{n-1}{2}\rfloor$}
           % \STATE (If $n$ odd, stop at $k =\lfloor\frac{n-1}{2}\rfloor-1 )$
            \STATE Compute \begin{align*}
                &\bigoplus_{\rho\in\rho_1\otimes\rho_{k-1}}\mathcal{F}(\Theta)_\rho =\left(C_{\rho_1,\rho_{k-1}}^{\dagger}\left[\mathcal{F}_{\rho_1}\otimes\mathcal{F}_{\rho_{k-1}}\right]^{-1}\beta_{\rho_1,\rho_2}C_{\rho_1,\rho_{k-1}}\right)^\dagger.
            \end{align*}
        \ENDFOR
        \STATE \textbf{Return} $\mathcal{F}(\Theta)$ (up to group action).
    \end{algorithmic}
\end{algorithm}
%\subsubsection{General algorithm for finite non-commutative groups}
%The three cases that we have just studied suggest a heuristic procedure to invert the Bispectrum. Given a group $G$ and a complete family of irreps $\{\rho_k\}_k$ where $\rho_0$ is the trivial representation. $\mathcal{F}(\Theta)_{\rho_0}$ can always be recovered from $\beta(\Theta)_{\rho_0,\rho_0}$. Then one can compute $\mathcal{F}(\Theta)_{\rho_0}$
\subsection{The 1D irreps of $D_n$}\label{app:1dirreps}
We recall the definition of the dihedral group $D_n$ given in~\eqref{def:dihedral}. The 1D irreps of $D_n$ can be found, e.g., in \citep{steinberg2011representation}. They are given by:
\begin{itemize}
    \item $\rho_{0}(g)=1$ for all $g\in D_n$.
    \item $\rho_{01}(g)=\begin{cases}
        1 \text{ if } g\in \langle a\rangle,\\
        -1 \text{ otherwise.}
    \end{cases}$
    \item If $n$ even, $\rho_{02}(g)=\begin{cases}
        1 \text{ if } g\in \langle a^2,\ x\rangle,\\
        -1 \text{ otherwise.}
    \end{cases}$
    \item If $n$ even, $\rho_{03}(g)=\begin{cases}
        1 \text{ if } g\in \langle a^2,\ ax\rangle,\\
        -1 \text{ otherwise.}
    \end{cases}$
\end{itemize}
To exemplify the 1D irreps, we give their values for $D_4$ in Table~\ref{tab:1dirreps}.
\begin{table}[H]
    \centering
    \begin{tabular}{|c||c|c|c|c|c|c|c|c|}
        \hline
         $g$&$e$&$a$&$a^2$&$a^3$&$x$&$ax$&$a^2x$&$a^3x$ \\
         \hline
         \hline
         $\rho_{0}(g)$&1&1&1&1&1&1&1&1 \\
         \hline
         $\rho_{01}(g)$&1&1&1&1&-1&-1&-1&-1 \\
         \hline
         $\rho_{02}(g)$&1&-1&1&-1&1&-1&1&-1 \\
         \hline
         $\rho_{03}(g)$&1&-1&1&-1&-1&1&-1&1 \\
         \hline
    \end{tabular}
    \caption{1D irreps of $D_4$.}
    \label{tab:1dirreps}
\end{table}
\subsubsection{Generation of the coefficients of $D_n$}\label{app:generating_dihedral}
Theorem~\ref{thm:dihedral_inversion} makes two assertions that we verify explicitly in this appendix. First, it is said that $\rho_{k}$ is in the tensor decomposition \eqref{eqn:non-commutative-bs} (TD) of $\rho_1\otimes\rho_{k-1}$ for $k = 2,...,\lfloor\frac{n-1}{2}\rfloor$ so that the iteration of Algorithm~\ref{alg:dihedral_alg} recovers all the Fourier coefficients associated to the 2D irreps. The second assertion to verify is that $\rho_{01}$ is in the TD of $ \rho_1\otimes\rho_1$ and, for $n$ even,  $\rho_{02},\rho_{03}$ in the TD of $\rho_1\otimes\rho_{\frac{n}{2}-1}$ so that the Fourier coefficients associated to the 1D irreps are also recovered, asserting the validity of the inversion procedure. We provide an analytical proof of these two assertions. The proof is based on the theory of \emph{character functions}.

\begin{definition}{\citep{steinberg2011representation}}
    Given a group $G$ and a representation $\rho$, the \emph{character} of $\rho$ is the function $\chi_\rho:G\rightarrow\mathbb{R}:g\mapsto\mathrm{Tr}(\rho(g))$. $\chi_\rho$ is said to be an \emph{irreducible} character if $\rho$ is an irreducible representation.
\end{definition}
The character function $\chi_\rho$ is a \emph{class function} on $G$, i.e., $\chi_\rho$ is constant on a conjugacy class of $G$. The space of class functions on a finite group $G$, written $\mathcal{S}_G$, can be equipped with an inner product $\langle\cdot,\cdot\rangle_G:\mathcal{S}_G\times \mathcal{S}_G\mapsto \mathbb{C}$ such that for $u,v\in \mathcal{S}_G$, we have
\begin{equation}
    \langle u, v\rangle_G:=\frac{1}{|G|}\sum_{g\in|G|}u(g)\overline{v(g)}.
\end{equation}
The irreducible characters form an orthonormal basis w.r.t  $\langle\cdot,\cdot\rangle_G$ for $\mathcal{S}_G$ \citep{steinberg2011representation}, i.e.,
\begin{equation}
    \langle \chi_{\rho_a},\chi_{\rho_b}\rangle_G=\begin{cases}
        1\text{ if } \rho_a\sim \rho_b,\\
        0 \text{ otherwise.}
    \end{cases}
\end{equation}
Therefore, for $\rho_a, \rho_b$ two irreps of $G$, $\rho_c$ is in the TD of $\rho_a\otimes\rho_b$ if and only if $ \langle \chi_{\rho_a\otimes \rho_b},\chi_{\rho_c}\rangle_G\neq 0$. Let us apply this to the 2D irreps of $D_n$ ($n>2$). Let $\rho_i,\rho_j,\rho_k$ be three irreps defined as in~\eqref{eq:2d_irreps_dn}.
Notice that we have $\mathcal{X}_{\rho}(a^lx)=0$. This yields
\begin{align*}
    \langle \mathcal{X}_{\rho_i\otimes\rho_j},\mathcal{X}_{\rho_k}\rangle &=\frac{1}{2n}\sum_{l=1}^n \mathcal{X}_{\rho_i\otimes\rho_j}(a^l)\mathcal{X}_{\rho_k}(a^l)\\
    &=\frac{1}{2n}\sum_{l=1}^n \mathcal{X}_{\rho_i}(a^l)\mathcal{X}_{\rho_j}(a^l)\mathcal{X}_{\rho_k}(a^l)\\
    &=\frac{1}{2n}\sum_{l=1}^n 8\cos(\omega_l i)\cos(\omega_l j)\cos(\omega_l k)\\
    %&=\frac{1}{2n}\sum_{l=1}^n 4\left(\cos(\omega_l (i+j))+\cos(\omega_l (j-i))\right)\cos(\omega_l k)\\
    &=\frac{1}{n}\sum_{l=1}^n \cos(\omega_l (i+j+k))+\cos(\omega_l (i+j-k))\\ &\ +\cos(\omega_l (j-i+k))+\cos(\omega_l (j-i-k)).
\end{align*}
Recall that $\sum_{l=1}^n\cos(\omega_l m)=\begin{cases}
        0 \text{ if } m\neq 0\\
        n \text{ if } m = 0
    \end{cases}$. Therefore, if we assume $0\leq i\leq j$ without loss of generality, $ \langle \mathcal{X}_{\rho_i\otimes\rho_j},\mathcal{X}_{\rho_k}\rangle\neq 0$ if and only if
\begin{equation*}
    \begin{cases}
        k = i+j \text{ or}\\
        k = j-i.
    \end{cases}.
\end{equation*}
Therefore, based on Definition~\ref{def:Bispectrum}, by utilizing $\beta_{\rho_i,\rho_j}$, $\mathcal{F}(\Theta)_{\rho_i}$ and $\mathcal{F}(\Theta)_{\rho_j}$, we can compute  $\mathcal{F}(\Theta)_{\rho_k}$ for $k\in\{i+j,\ j-i\}$. By iterating, if $\mathcal{F}(\Theta)_{\rho_1}$ is known,  $\beta_{\rho_1,\rho_1}$ gives $\mathcal{F}(\Theta)_{\rho_2}$. Then, $\beta_{\rho_1,\rho_2}$ can be leveraged to obtain  $\mathcal{F}(\Theta)_{\rho_3}$. Continuing the procedure provides $\mathcal{F}(\Theta)_{\rho_k}$ for $k=2,...,\lfloor \frac{n-1}{2}\rfloor$ by using $\beta_{\rho_1,\rho_{k-1}}$. We have thus obtained the Fourier coefficients associated to all the 2D irreps.

It remains to show that the iteration also recovered the Fourier coefficients associated to the 1D irreps. This is because $\rho_{01}$ is in the TD of $\rho_1\otimes\rho_1$ and, for $n$ even, $\rho_{02}, \rho_{03}$ are in the TD of $\rho_1\otimes\rho_{\frac{n}{2}-1}$. Indeed,
\begin{align*}
    \langle \mathcal{X}_{\rho_1\otimes\rho_1},\mathcal{X}_{\rho_{01}}\rangle &=\frac{1}{2n}\sum_{l=1}^n \mathcal{X}_{\rho_1\otimes\rho_1}(a^l)\mathcal{X}_{\rho_{01}}(a^l)\\
    &=\frac{1}{2n}\sum_{l=1}^n 4\cos(\omega_l)^2\\
    &=\frac{1}{n}\sum_{l=1}^n 1+\cos(2\omega_l)=1.
\end{align*}
Moreover, for $n$ even and $\rho\in\{\rho_{02}, \rho_{03}\}$, we have
\begin{align*}
    \langle \mathcal{X}_{\rho_1\otimes\rho_{\frac{n}{2}-1}},\mathcal{X}_\rho\rangle &=\frac{1}{2n}\sum_{l=1}^n \mathcal{X}_{\rho_1\otimes\rho_{\frac{n}{2}-1}}(a^l)\mathcal{X}_\rho(a^l)\\
    &=\frac{1}{n}\sum_{l=1}^n 1+\cos\left(\omega_l\left(\frac{n}{2}-2\right)\right)(-1)^l\\
    &=1.
\end{align*}
In conclusion, the procedure of Algorithm \ref{alg:dihedral_alg} recovers all the Fourier coefficients and the selective $G$-Bispectrum.


\subsection{The Clebsch-Gordan matrices on $D_n$}\label{app:clebsch_gordan}
The matrix algebra properties that we use in this subsection can be found, e.g., in \citep{bhatia97}. Recall from Theorem~\ref{thm:commutative_Bispectrum} the (implicit) definition of the Clebsch-Gordan matrices:
\begin{equation}\label{eq:cleb_gordan}
    (\rho_1\otimes\rho_2)(g) = C_{\rho_1,\rho_2}\left[ \bigoplus_{\rho\in\mathcal{R}} \rho(g) \right]C_{\rho_1,\rho_2}^\dagger,
\end{equation}
where $C_{\rho_1,\rho_2}^\dagger C_{\rho_1,\rho_2}=I$.
We only consider the case of $\rho_1,\rho_2$ both 2d irreps of $D_n$ since otherwise, the Clebsch-Gordan matrix is the scalar $1$.
First notice that $(\rho_1\otimes\rho_2)(g)$, is an orthonormal matrix. Indeed, using the properties of the Kronecker product, we obtain (``$(g)$'' omitted for clarity):
\begin{align*}
    (\rho_1\otimes\rho_2)(\rho_1\otimes\rho_2)^\dagger &= (\rho_1\otimes\rho_2)(\rho_1^\dagger\otimes\rho_2^\dagger)\\
    &=(\rho_1\rho_1^\dagger)\otimes (\rho_2\rho_2^\dagger)\\
    &=I \otimes I = I
\end{align*}
For $Q$ a real orthogonal matrix ($Q^T Q = I$), and $VSV^T$ with $V^T V = I$, a real Schur decomposition of $Q$, it is known that $S$ is block diagonal with blocks of size $1\times 1$ or $2\times 2$. These blocks are themselves orthogonal matrices. Therefore, the real Schur decomposition is the decomposition in~\eqref{eq:cleb_gordan} up to permutations. In order for $S$ to represent exactly the irreps from~\eqref{eq:2d_irreps_dn}, the non-zero sub-diagonal elements should all be positive. If not, the symmetric element is positive and a permutation $P$ must be added to exchange their positions: $Q = (VP)(P^TSP)(VP)^T$. $P$ permutes the two columns of $V$ associated with the permuted $2\times 2$ block of $S$.
\begin{algorithm}
    \caption{Compute Clebsch-Gordan matrices on $D_n$}
    \begin{algorithmic}[1]
        \STATE \textbf{Input:} $\rho_1, \rho_2$, two 2d irreps of $D_n$.
        \STATE Pick any $g\in D_n$, e.g., $g = a$.
        \STATE Compute  $(\rho_1\otimes\rho_2)(g) = (VP)(P^TSP)(VP)^T$, a valid real Schur form.
        \STATE Set $C_{\rho_1,\rho_2} = VP$.
        \STATE Set $\bigoplus_{\rho\in\mathcal{R}} \rho(g) = P^TSP$
        \STATE \textbf{Return:} $C_{\rho_1,\rho_2}, \bigoplus_{\rho\in\mathcal{R}}\rho(g)$.
    \end{algorithmic}
\end{algorithm}




\section{Bispectrum inversion for octahedral and full octahedral groups}\label{app:octahedral}
We provide a sketch of the procedure to retrieve $\mathcal{F}(\Theta)$ given $\beta(\Theta)$ for the \href{https://en.wikipedia.org/wiki/Octahedral_symmetry}{octahedral group} and the \href{https://en.wikiversity.org/wiki/Full_octahedral_group}{full octahedral group}. These two groups are available in the \href{https://quva-lab.github.io/escnn/api/escnn.group.html}{\texttt{escnn} library}. These groups are easier to deal with than the cyclic and dihedral groups presented in the paper, given that they do not come from a \textit{family} of groups. Indeed, our proofs for the cyclic (resp. dihedral) groups needed to work for \textit{all} cyclic groups $C_N$, and for all dihedral groups $D_N$, for all $N$. The octahedral and full octahedral groups are only two groups.

\subsection{Octahedral group}
\begin{table}
        \centering
        \begin{tabular}{|c||c|c|c|c|c|}
            \hline
            $\otimes$ &$\rho_0$&$\rho_1$&$\rho_2$&$\rho_3$&$\rho_4$\\
            \hline
            \hline
            $\rho_0$&10000 & 01000 & 00100 & 00010 & 00001 \\
            \hline
            $\rho_1$&01000 & 11110 & 01111 & 01100 & 00100 \\
            \hline
            $\rho_2$&00100 & 01111 & 11110 & 01100 & 01000 \\
            \hline
            $\rho_3$&00010 & 01100 & 01100 & 10011 & 00010 \\
            \hline
            $\rho_4$ &00001 & 00100 & 01000 & 00010 & 10000 \\
            \hline
        \end{tabular}
        \caption{Kronecker table of the octahedral group using \texttt{escnn}. For the binary word at position $i,j$ in the table, the $k$th `letter' is $1$ if $\rho_k\in\rho_i\otimes \rho_j$, $0$ otherwise.}
        \label{tab:octa}
    \end{table}

The octahedral group has 24 elements and 5 irreps. We can compute its Kronecker table, either manually using characters $\chi$ or using a Python package such as \texttt{escnn}. We give its Kronecker table below (Table~\ref{tab:octa}), where each column/row represents one irrep, labelled $\rho_0, \rho_1, ..., \rho_4$.

We apply the procedure from Algorithms~\ref{alg:cyclic_alg}, \ref{alg:commutative_alg} and \ref{alg:dihedral_alg} to Table~\ref{tab:octa}. This procedure relies on the use of Theorem~\ref{def:Bispectrum}. We first select the bispectral coefficient $\beta_{\rho_0, \rho_0}$ (we omit ``$(\Theta)$'' for clarity) to get the component $\mathcal{F}_{\rho_0 }$ where $\rho_0$ is the trivial representation. Next, we choose $\beta_{\rho_0,\rho_1}$ and use $\mathcal{F}_{\rho_0}$  to obtain $\mathcal{F}_{\rho_1}$ (we know from Table~\ref{tab:octa} that $\rho_1\in\rho_0\otimes\rho_1$) up to an indeterminacy which is a transformation in $\mathrm{O}(3)$ and corresponds to the indeterminacy factor from Appendix~\ref{sec:indeterminacy}. Then, we select  $\beta_{\rho_1, \rho_1}$ to get the Fourier components $\mathcal{F}_{\rho_2}, \mathcal{F}_{\rho_3}$. Lastly, we select $\beta_{\rho_1, \rho_2}$ to get the missing Fourier component $\mathcal{F}_{\rho_4}$.

In summary, we only need $4$ bispectral coefficients ($\beta_{\rho_0, \rho_0}, \beta_{\rho_1, \rho_0}, \beta_{\rho_1, \rho_1}, \beta_{\rho_1, \rho_2}$) instead of $5^2 = 25$ in order to get the five Fourier components, i.e., the full Fourier transform of the signal. In total, this involves $1+9+81+81=172$ scalar coefficients.
\subsection{Full octahedral group}
\begin{table*}[h]
    \centering
    \minuscule
    \begin{tabular}{|c||c|c|c|c|c|c|c|c|c|c|c|}
    \hline
            $\otimes$ &$\rho_0$&$\rho_1$&$\rho_2$&$\rho_3$&$\rho_4$&$\rho_5$&$\rho_6$&$\rho_7$&$\rho_8$&$\rho_9$\\
    \hline
    $\rho_0$&1000000000 & 0100000000 & 0010000000 & 0001000000 & 0000100000 & 0000010000 & 0000001000 & 0000000100 & 0000000010 & 0000000001 \\
    \hline
    $\rho_1$&0100000000 & 1111000000 & 0111100000 & 0110000000 & 0010000000 & 0000001000 & 0000011110 & 0000001111 & 0000001100 & 0000000100 \\
    \hline
    $\rho_2$&0010000000 & 0111100000 & 1111000000 & 0110000000 & 0100000000 & 0000000100 & 0000001111 & 0000011110 & 0000001100 & 0000001000 \\
    \hline
    $\rho_3$&0001000000 & 0110000000 & 0110000000 & 1001100000 & 0001000000 & 0000000010 & 0000001100 & 0000001100 & 0000010011 & 0000000010 \\
    \hline
    $\rho_4$&0000100000 & 0010000000 & 0100000000 & 0001000000 & 1000000000 & 0000000001 & 0000000100 & 0000001000 & 0000000010 & 0000010000 \\
    \hline
    $\rho_5$&0000010000 & 0000001000 & 0000000100 & 0000000010 & 0000000001 & 1000000000 & 0100000000 & 0010000000 & 0001000000 & 0000100000 \\
    \hline
    $\rho_6$&0000001000 & 0000011110 & 0000001111 & 0000001100 & 0000000100 & 0100000000 & 1111000000 & 0111100000 & 0110000000 & 0010000000 \\
    \hline
    $\rho_7$&0000000100 & 0000001111 & 0000011110 & 0000001100 & 0000001000 & 0010000000 & 0111100000 & 1111000000 & 0110000000 & 0100000000 \\
    \hline
    $\rho_8$&0000000010 & 0000001100 & 0000001100 & 0000010011 & 0000000010 & 0001000000 & 0110000000 & 0110000000 & 1001100000 & 0001000000 \\
    \hline
    $\rho_9$&0000000001 & 0000000100 & 0000001000 & 0000000010 & 0000010000 & 0000100000 & 0010000000 & 0100000000 & 0001000000 & 1000000000\\
    \hline
    \end{tabular}
    \caption{Kronecker table of  full octahedral group using \texttt{escnn}. For the binary word at position $i,j$ in the table, the $k$th `letter' is $1$ if $\rho_k\in\rho_i\otimes \rho_j$, $0$ otherwise.}
    \label{tab:full_octa}
\end{table*}
The full octahedral group has $48$ elements and $10$ irreps. Again, we can compute its Kronecker table using a Python package such as \texttt{escnn}. We give its Kronecker table below (Table~\ref{tab:full_octa}), where each column/row represents one irrep, labelled $\rho_0, \rho_1, ..., \rho_{9}$.

We apply the procedure from Algorithms~\ref{alg:cyclic_alg}, \ref{alg:commutative_alg} and \ref{alg:dihedral_alg} to Table~\ref{tab:full_octa}. Again, this procedure relies on the use of Theorem~\ref{def:Bispectrum}. $\beta_{\rho_0,\rho_0}$ (we omit ``$(\Theta)$'' for clarity) allows to compute $\mathcal{F}_{\rho_0}$ directly, such as in Algorithm~\ref{alg:dihedral_alg}. Then, from $\beta_{\rho_0,\rho_6}$, we obtain $\mathcal{F}_{\rho_6}$ up to an unknown group action in $\mathrm{O}(3)$. Then, using $\beta_{\rho_6,\rho_6}$ and $\mathcal{F}_{\rho_6}$, we obtain $\mathcal{F}_{\rho_1}$,$\mathcal{F}_{\rho_2}$ and $\mathcal{F}_{\rho_ 3}$. Next, leveraging $\beta_{\rho_1,\rho_2}$, $\mathcal{F}_{\rho_1}$ and $\mathcal{F}_{\rho_2}$, we obtain  $\mathcal{F}_{\rho_4}$. Using $\beta_{\rho_1,\rho_6}$, $\mathcal{F}_{\rho_1}$ and $\mathcal{F}_{\rho_6}$, we obtain  $\mathcal{F}_{\rho_5}$, $\mathcal{F}_{\rho_7}$, $\mathcal{F}_{\rho_8}$. Finally, with $\beta_{\rho_1,\rho_7}$, $\mathcal{F}_{\rho_1}$ and $\mathcal{F}_{\rho_7}$, we obtain the last coefficient $\mathcal{F}_{\rho_9}$. Hence we have recovered all the Fourier coefficients using only $\beta_{\rho_0,\rho_0}$, $\beta_{\rho_0,\rho_6}$, $\beta_{\rho_6,\rho_6}$, $\beta_{\rho_1,\rho_2}$, $\beta_{\rho_1,\rho_6}$, $\beta_{\rho_1,\rho_7}$, thus a total of $6$ bispectral coefficients instead of $100$. In terms of scalar values, this involves $1+9+81+81+81+81=334$ coefficients.


\section{Training of the $G$-CNN architecture}\label{app:training}
The $\mathrm{SO}(2)$-MNIST/EMNIST datasets are obtained after applying random planar rotations on each image of the datasets MNIST \citep{lecun2010mnist}, EMNIST \citep{cohen2017emnist} respectively. In the case of $\mathrm{O}(2)$-MNIST/EMNIST, in addition to a planar rotation, a reflection is applied with probability $\frac{1}{2}$.  The original size of each image is conserved. The size of the training sets are $60\ 000$ and $88\ 800$ for MNIST and EMNIST, respectively. %The size of the test sets are $10\ 000$ and $14\ 800$. For our experiments, $20\%$ of the training datasets are used for validation.
%\subsection{$C_8$-CNNs on $\mathrm{SO}(2)$-MNIST/EMNIST/CIFAR-10}

We conserve the architecture of \citep{sanborn2023general}. For all invariant layers, being the $G$-TC, the selective $G$-Bispectrum and the Avg/Max $G$-pooling, the architecture is composed of a $C_8/D_8$-convolutional block with $K$ filters (see Table~\ref{tab:classification}). Then, the invariant layer is applied before feeding the output to a MLP. The MLP is composed of $3$ fully-connected layers with ReLU non-linearity. A final fully-connected linear layer is applied for classification. The vector of the output sizes of these layers is given by $[o1,o2,o3,ol]$ respectively. $o2=o3 = 64$ and $ol$ is equal to the number of classes of the dataset. $o1$ is tuned to reach the parameter count from Table~\ref{tab:classification}.
\begin{comment}
\newpage

\section*{NeurIPS Paper Checklist}
\begin{enumerate}

\item {\bf Claims}
    \item[] Question: Do the main claims made in the abstract and introduction accurately reflect the paper's contributions and scope?
    \item[] Answer: \answerYes{} % Replace by \answerYes{}, \answerNo{}, or \answerNA{}.
    \item[] Justification: The claims that we make in the abstract, i.e., theoretical and practical guarantees of the $G$-Bispectrum are respectively  proven and studied in Sections \ref{sec:completeness} and \ref{sec:experiments}.%\justificationTODO{}
    \item[] Guidelines:
    \begin{itemize}
        \item The answer NA means that the abstract and introduction do not include the claims made in the paper.
        \item The abstract and/or introduction should clearly state the claims made, including the contributions made in the paper and important assumptions and limitations. A No or NA answer to this question will not be perceived well by the reviewers.
        \item The claims made should match theoretical and experimental results, and reflect how much the results can be expected to generalize to other settings.
        \item It is fine to include aspirational goals as motivation as long as it is clear that these goals are not attained by the paper.
    \end{itemize}

\item {\bf Limitations}
    \item[] Question: Does the paper discuss the limitations of the work performed by the authors?
    \item[] Answer: \answerYes{} % Replace by \answerYes{}, \answerNo{}, or \answerNA{}.
    \item[] Justification: In Section~\ref{sec:experiments}, we discuss the pros and cons of the selective $G$-Bispectrum compared to the other possible invariant layers. %\justificationTODO{}
    \item[] Guidelines:
    \begin{itemize}
        \item The answer NA means that the paper has no limitation while the answer No means that the paper has limitations, but those are not discussed in the paper.
        \item The authors are encouraged to create a separate "Limitations" section in their paper.
        \item The paper should point out any strong assumptions and how robust the results are to violations of these assumptions (e.g., independence assumptions, noiseless settings, model well-specification, asymptotic approximations only holding locally). The authors should reflect on how these assumptions might be violated in practice and what the implications would be.
        \item The authors should reflect on the scope of the claims made, e.g., if the approach was only tested on a few datasets or with a few runs. In general, empirical results often depend on implicit assumptions, which should be articulated.
        \item The authors should reflect on the factors that influence the performance of the approach. For example, a facial recognition algorithm may perform poorly when image resolution is low or images are taken in low lighting. Or a speech-to-text system might not be used reliably to provide closed captions for online lectures because it fails to handle technical jargon.
        \item The authors should discuss the computational efficiency of the proposed algorithms and how they scale with dataset size.
        \item If applicable, the authors should discuss possible limitations of their approach to address problems of privacy and fairness.
        \item While the authors might fear that complete honesty about limitations might be used by reviewers as grounds for rejection, a worse outcome might be that reviewers discover limitations that aren't acknowledged in the paper. The authors should use their best judgment and recognize that individual actions in favor of transparency play an important role in developing norms that preserve the integrity of the community. Reviewers will be specifically instructed to not penalize honesty concerning limitations.
    \end{itemize}

\item {\bf Theory Assumptions and Proofs}
    \item[] Question: For each theoretical result, does the paper provide the full set of assumptions and a complete (and correct) proof?
    \item[] Answer: \answerYes{} % Replace by \answerYes{}, \answerNo{}, or \answerNA{}.
    \item[] Justification: The previously known theorems and results are referenced and each new theorem is accompanied by a proof, usually proposed in the appendices. %\justificationTODO{}
    \item[] Guidelines:
    \begin{itemize}
        \item The answer NA means that the paper does not include theoretical results.
        \item All the theorems, formulas, and proofs in the paper should be numbered and cross-referenced.
        \item All assumptions should be clearly stated or referenced in the statement of any theorems.
        \item The proofs can either appear in the main paper or the supplemental material, but if they appear in the supplemental material, the authors are encouraged to provide a short proof sketch to provide intuition.
        \item Inversely, any informal proof provided in the core of the paper should be complemented by formal proofs provided in appendix or supplemental material.
        \item Theorems and Lemmas that the proof relies upon should be properly referenced.
    \end{itemize}

    \item {\bf Experimental Result Reproducibility}
    \item[] Question: Does the paper fully disclose all the information needed to reproduce the main experimental results of the paper to the extent that it affects the main claims and/or conclusions of the paper (regardless of whether the code and data are provided or not)?
    \item[] Answer: \answerYes{} % Replace by \answerYes{}, \answerNo{}, or \answerNA{}.
    \item[] Justification: We provide all the details about the architectures of the $G$-CNNs that we use in Section \ref{sec:experiments} and Appendix~\ref{app:training}. The code will be made available upon publication, with precise guidance to reproduce the results.
    \item[] Guidelines:
    \begin{itemize}
        \item The answer NA means that the paper does not include experiments.
        \item If the paper includes experiments, a No answer to this question will not be perceived well by the reviewers: Making the paper reproducible is important, regardless of whether the code and data are provided or not.
        \item If the contribution is a dataset and/or model, the authors should describe the steps taken to make their results reproducible or verifiable.
        \item Depending on the contribution, reproducibility can be accomplished in various ways. For example, if the contribution is a novel architecture, describing the architecture fully might suffice, or if the contribution is a specific model and empirical evaluation, it may be necessary to either make it possible for others to replicate the model with the same dataset, or provide access to the model. In general. releasing code and data is often one good way to accomplish this, but reproducibility can also be provided via detailed instructions for how to replicate the results, access to a hosted model (e.g., in the case of a large language model), releasing of a model checkpoint, or other means that are appropriate to the research performed.
        \item While NeurIPS does not require releasing code, the conference does require all submissions to provide some reasonable avenue for reproducibility, which may depend on the nature of the contribution. For example
        \begin{enumerate}
            \item If the contribution is primarily a new algorithm, the paper should make it clear how to reproduce that algorithm.
            \item If the contribution is primarily a new model architecture, the paper should describe the architecture clearly and fully.
            \item If the contribution is a new model (e.g., a large language model), then there should either be a way to access this model for reproducing the results or a way to reproduce the model (e.g., with an open-source dataset or instructions for how to construct the dataset).
            \item We recognize that reproducibility may be tricky in some cases, in which case authors are welcome to describe the particular way they provide for reproducibility. In the case of closed-source models, it may be that access to the model is limited in some way (e.g., to registered users), but it should be possible for other researchers to have some path to reproducing or verifying the results.
        \end{enumerate}
    \end{itemize}


\item {\bf Open access to data and code}
    \item[] Question: Does the paper provide open access to the data and code, with sufficient instructions to faithfully reproduce the main experimental results, as described in supplemental material?
    \item[] Answer: \answerYes{} % Replace by \answerYes{}, \answerNo{}, or \answerNA{}.
    \item[] Justification: The code will only be made available after publication to preserve the anonymous component of the review process. This code will include precise guidance to reproduce the results.%\justificationTODO{}
    \item[] Guidelines:
    \begin{itemize}
        \item The answer NA means that paper does not include experiments requiring code.
        \item Please see the NeurIPS code and data submission guidelines (\url{https://nips.cc/public/guides/CodeSubmissionPolicy}) for more details.
        \item While we encourage the release of code and data, we understand that this might not be possible, so “No” is an acceptable answer. Papers cannot be rejected simply for not including code, unless this is central to the contribution (e.g., for a new open-source benchmark).
        \item The instructions should contain the exact command and environment needed to run to reproduce the results. See the NeurIPS code and data submission guidelines (\url{https://nips.cc/public/guides/CodeSubmissionPolicy}) for more details.
        \item The authors should provide instructions on data access and preparation, including how to access the raw data, preprocessed data, intermediate data, and generated data, etc.
        \item The authors should provide scripts to reproduce all experimental results for the new proposed method and baselines. If only a subset of experiments are reproducible, they should state which ones are omitted from the script and why.
        \item At submission time, to preserve anonymity, the authors should release anonymized versions (if applicable).
        \item Providing as much information as possible in supplemental material (appended to the paper) is recommended, but including URLs to data and code is permitted.
    \end{itemize}


\item {\bf Experimental Setting/Details}
    \item[] Question: Does the paper specify all the training and test details (e.g., data splits, hyperparameters, how they were chosen, type of optimizer, etc.) necessary to understand the results?
    \item[] Answer: \answerYes{} % Replace by \answerYes{}, \answerNo{}, or \answerNA{}.
    \item[] Justification: We provide in Section~\ref{sec:experiments} the details about the datasets and the architectures of the models implemented for experiments. The optimization details and hyper-parameters choice will be made available with the code upon publication.%\justificationTODO{}
    \item[] Guidelines:
    \begin{itemize}
        \item The answer NA means that the paper does not include experiments.
        \item The experimental setting should be presented in the core of the paper to a level of detail that is necessary to appreciate the results and make sense of them.
        \item The full details can be provided either with the code, in appendix, or as supplemental material.
    \end{itemize}

\item {\bf Experiment Statistical Significance}
    \item[] Question: Does the paper report error bars suitably and correctly defined or other appropriate information about the statistical significance of the experiments?
    \item[] Answer: \answerYes{} % Replace by \answerYes{}, \answerNo{}, or \answerNA{}.
    \item[] Justification: All the experiments, i.e., Table~\ref{tab:classification} and Figures~\ref{fig:trainings} and \ref{fig:Kfilters} represent the standard deviations of the experiments performed.%\justificationTODO{}
    \item[] Guidelines:
    \begin{itemize}
        \item The answer NA means that the paper does not include experiments.
        \item The authors should answer "Yes" if the results are accompanied by error bars, confidence intervals, or statistical significance tests, at least for the experiments that support the main claims of the paper.
        \item The factors of variability that the error bars are capturing should be clearly stated (for example, train/test split, initialization, random drawing of some parameter, or overall run with given experimental conditions).
        \item The method for calculating the error bars should be explained (closed form formula, call to a library function, bootstrap, etc.)
        \item The assumptions made should be given (e.g., Normally distributed errors).
        \item It should be clear whether the error bar is the standard deviation or the standard error of the mean.
        \item It is OK to report 1-sigma error bars, but one should state it. The authors should preferably report a 2-sigma error bar than state that they have a 96\% CI, if the hypothesis of Normality of errors is not verified.
        \item For asymmetric distributions, the authors should be careful not to show in tables or figures symmetric error bars that would yield results that are out of range (e.g. negative error rates).
        \item If error bars are reported in tables or plots, The authors should explain in the text how they were calculated and reference the corresponding figures or tables in the text.
    \end{itemize}

\item {\bf Experiments Compute Resources}
    \item[] Question: For each experiment, does the paper provide sufficient information on the computer resources (type of compute workers, memory, time of execution) needed to reproduce the experiments?
    \item[] Answer: \answerYes{} % Replace by \answerYes{}, \answerNo{}, or \answerNA{}.
    \item[] Justification: We give in section \ref{sec:experiments} the running times of the algorithms and the hardware used for training. %\justificationTODO{}
    \item[] Guidelines:
    \begin{itemize}
        \item The answer NA means that the paper does not include experiments.
        \item The paper should indicate the type of compute workers CPU or GPU, internal cluster, or cloud provider, including relevant memory and storage.
        \item The paper should provide the amount of compute required for each of the individual experimental runs as well as estimate the total compute.
        \item The paper should disclose whether the full research project required more compute than the experiments reported in the paper (e.g., preliminary or failed experiments that didn't make it into the paper).
    \end{itemize}

\item {\bf Code Of Ethics}
    \item[] Question: Does the research conducted in the paper conform, in every respect, with the NeurIPS Code of Ethics \url{https://neurips.cc/public/EthicsGuidelines}?
    \item[] Answer: \answerYes{} % Replace by \answerYes{}, \answerNo{}, or \answerNA{}.
    \item[] Justification: This paper is not involved with specific ethical issues, except that it contributes to the field of image processing. The latter field is related to the ethical issue of face-recognition and automatic surveillance.%\justificationTODO{}
    \item[] Guidelines:
    \begin{itemize}
        \item The answer NA means that the authors have not reviewed the NeurIPS Code of Ethics.
        \item If the authors answer No, they should explain the special circumstances that require a deviation from the Code of Ethics.
        \item The authors should make sure to preserve anonymity (e.g., if there is a special consideration due to laws or regulations in their jurisdiction).
    \end{itemize}


\item {\bf Broader Impacts}
    \item[] Question: Does the paper discuss both potential positive societal impacts and negative societal impacts of the work performed?
    \item[] Answer: \answerNA{} % Replace by \answerYes{}, \answerNo{}, or \answerNA{}.
    \item[] Justification: This work is a research work to propose a new layer in the architecture of $G$-CNNS. It has no direct societal impacts beyond the fact that it contributes to image-processing, which can be used for malicious purposes. %\justificationTODO{}
    \item[] Guidelines:
    \begin{itemize}
        \item The answer NA means that there is no societal impact of the work performed.
        \item If the authors answer NA or No, they should explain why their work has no societal impact or why the paper does not address societal impact.
        \item Examples of negative societal impacts include potential malicious or unintended uses (e.g., disinformation, generating fake profiles, surveillance), fairness considerations (e.g., deployment of technologies that could make decisions that unfairly impact specific groups), privacy considerations, and security considerations.
        \item The conference expects that many papers will be foundational research and not tied to particular applications, let alone deployments. However, if there is a direct path to any negative applications, the authors should point it out. For example, it is legitimate to point out that an improvement in the quality of generative models could be used to generate deepfakes for disinformation. On the other hand, it is not needed to point out that a generic algorithm for optimizing neural networks could enable people to train models that generate Deepfakes faster.
        \item The authors should consider possible harms that could arise when the technology is being used as intended and functioning correctly, harms that could arise when the technology is being used as intended but gives incorrect results, and harms following from (intentional or unintentional) misuse of the technology.
        \item If there are negative societal impacts, the authors could also discuss possible mitigation strategies (e.g., gated release of models, providing defenses in addition to attacks, mechanisms for monitoring misuse, mechanisms to monitor how a system learns from feedback over time, improving the efficiency and accessibility of ML).
    \end{itemize}

\item {\bf Safeguards}
    \item[] Question: Does the paper describe safeguards that have been put in place for responsible release of data or models that have a high risk for misuse (e.g., pretrained language models, image generators, or scraped datasets)?
    \item[] Answer: \answerNA{} % Replace by \answerYes{}, \answerNo{}, or \answerNA{}.
    \item[] Justification: This paper does not leverage models with high risk of misuse.%\justificationTODO{}
    \item[] Guidelines:
    \begin{itemize}
        \item The answer NA means that the paper poses no such risks.
        \item Released models that have a high risk for misuse or dual-use should be released with necessary safeguards to allow for controlled use of the model, for example by requiring that users adhere to usage guidelines or restrictions to access the model or implementing safety filters.
        \item Datasets that have been scraped from the Internet could pose safety risks. The authors should describe how they avoided releasing unsafe images.
        \item We recognize that providing effective safeguards is challenging, and many papers do not require this, but we encourage authors to take this into account and make a best faith effort.
    \end{itemize}

\item {\bf Licenses for existing assets}
    \item[] Question: Are the creators or original owners of assets (e.g., code, data, models), used in the paper, properly credited and are the license and terms of use explicitly mentioned and properly respected?
    \item[] Answer: \answerYes{} % Replace by \answerYes{}, \answerNo{}, or \answerNA{}.
    \item[] Justification: We use the datasets MNIST \citep{deng2012mnist} and EMNIST \citep{cohen2017emnist}, which we properly mention in the main text. The URLs are displayed with hyper-references.%\justificationTODO{}
    \item[] Guidelines:
    \begin{itemize}
        \item The answer NA means that the paper does not use existing assets.
        \item The authors should cite the original paper that produced the code package or dataset.
        \item The authors should state which version of the asset is used and, if possible, include a URL.
        \item The name of the license (e.g., CC-BY 4.0) should be included for each asset.
        \item For scraped data from a particular source (e.g., website), the copyright and terms of service of that source should be provided.
        \item If assets are released, the license, copyright information, and terms of use in the package should be provided. For popular datasets, \url{paperswithcode.com/datasets} has curated licenses for some datasets. Their licensing guide can help determine the license of a dataset.
        \item For existing datasets that are re-packaged, both the original license and the license of the derived asset (if it has changed) should be provided.
        \item If this information is not available online, the authors are encouraged to reach out to the asset's creators.
    \end{itemize}

\item {\bf New Assets}
    \item[] Question: Are new assets introduced in the paper well documented and is the documentation provided alongside the assets?
    \item[] Answer: \answerNA{} % Replace by \answerYes{}, \answerNo{}, or \answerNA{}.
    \item[] Justification: We do not use new assets and rely on existing datasets.
    \item[] Guidelines:
    \begin{itemize}
        \item The answer NA means that the paper does not release new assets.
        \item Researchers should communicate the details of the dataset/code/model as part of their submissions via structured templates. This includes details about training, license, limitations, etc.
        \item The paper should discuss whether and how consent was obtained from people whose asset is used.
        \item At submission time, remember to anonymize your assets (if applicable). You can either create an anonymized URL or include an anonymized zip file.
    \end{itemize}

\item {\bf Crowdsourcing and Research with Human Subjects}
    \item[] Question: For crowdsourcing experiments and research with human subjects, does the paper include the full text of instructions given to participants and screenshots, if applicable, as well as details about compensation (if any)?
    \item[] Answer: \answerNA{} % Replace by \answerYes{}, \answerNo{}, or \answerNA{}.
    \item[] Justification: No participants were required during the experiments. %\justificationTODO{}
    \item[] Guidelines:
    \begin{itemize}
        \item The answer NA means that the paper does not involve crowdsourcing nor research with human subjects.
        \item Including this information in the supplemental material is fine, but if the main contribution of the paper involves human subjects, then as much detail as possible should be included in the main paper.
        \item According to the NeurIPS Code of Ethics, workers involved in data collection, curation, or other labor should be paid at least the minimum wage in the country of the data collector.
    \end{itemize}

\item {\bf Institutional Review Board (IRB) Approvals or Equivalent for Research with Human Subjects}
    \item[] Question: Does the paper describe potential risks incurred by study participants, whether such risks were disclosed to the subjects, and whether Institutional Review Board (IRB) approvals (or an equivalent approval/review based on the requirements of your country or institution) were obtained?
    \item[] Answer: \answerNA{} % Replace by \answerYes{}, \answerNo{}, or \answerNA{}.
    \item[] Justification: No human subjects were necessary to prepare this manuscript.%\justificationTODO{}
    \item[] Guidelines:
    \begin{itemize}
        \item The answer NA means that the paper does not involve crowdsourcing nor research with human subjects.
        \item Depending on the country in which research is conducted, IRB approval (or equivalent) may be required for any human subjects research. If you obtained IRB approval, you should clearly state this in the paper.
        \item We recognize that the procedures for this may vary significantly between institutions and locations, and we expect authors to adhere to the NeurIPS Code of Ethics and the guidelines for their institution.
        \item For initial submissions, do not include any information that would break anonymity (if applicable), such as the institution conducting the review.
    \end{itemize}

\end{enumerate}
\end{comment}

\end{document}
